% papers/wording/expected.tex                                      -*-LaTeX-*-
% SPDX-License-Identifier: Apache-2.0 WITH LLVM-exception
%
% Generated from the annotated headers in include/beman/expected/ via specgen.
% Do not edit by hand: change a header's //! docblocks and run `make wording`
% instead. This comment comes from papers/wording/preamble.tex, which is the
% one part of this file a person writes.
%
% --base-section-depth 2 puts each \rSec marker at the level the draft
% writes it at: [expected.unexpected] and its siblings are \rSec2, their
% subclauses \rSec3, numbering 22.8.3, 22.8.3.1, ... just as in
% source/utilities.tex. This file is therefore everything that sits
% *inside* the existing \rSec1[expected]{Expected objects} in the draft's
% source/utilities.tex, in clause order, ready to replace the current
% [expected.unexpected] through [expected.void] subclauses and add the new
% [expected.ref] one after them -- no \rSec1[expected] wrapper and no
% \input directives. [expected.general] and [expected.syn] are prose, not
% generated from any one declaration; see papers/expected-new.tex.

\rSec2[expected.unexpected]{Class template unexpected}

\rSec3[expected.un.general]{General}

\begin{codeblock}
template<class E>
class @\libglobal{unexpected}@ {
  // [expected.un.general] para 2: ill-formed instantiations

public:
  constexpr unexpected(const unexpected&) = default;
  constexpr unexpected(unexpected&&) = default;

  template<class Err = E>
    requires(!is_same_v<remove_cvref_t<Err>, unexpected> &&
             !is_same_v<remove_cvref_t<Err>, in_place_t> && is_constructible_v<E, Err>)
  constexpr explicit unexpected(Err&& e) noexcept(is_nothrow_constructible_v<E, Err>);

  template<class... Args>
    requires is_constructible_v<E, Args...>
  constexpr explicit unexpected(in_place_t, Args&&... args) noexcept(
      is_nothrow_constructible_v<E, Args...>);

  template<class U, class... Args>
    requires is_constructible_v<E, initializer_list<U>&, Args...>
  constexpr explicit unexpected(
      in_place_t, initializer_list<U> il,
      Args&&... args) noexcept(is_nothrow_constructible_v<E, initializer_list<U>&,
                                                          Args...>);

  constexpr unexpected& operator=(const unexpected&) = default;
  constexpr unexpected& operator=(unexpected&&) = default;

  constexpr const E& error() const& noexcept;
  constexpr E& error() & noexcept;
  constexpr const E&& error() const&& noexcept;
  constexpr E&& error() && noexcept;

  constexpr void swap(unexpected& other) noexcept(is_nothrow_swappable_v<E>);

  template<class E2>
  friend constexpr bool operator==(const unexpected& x, const unexpected<E2>& y);

  friend constexpr void swap(unexpected& x, unexpected& y) noexcept(noexcept(x.swap(y)))
    requires is_swappable_v<E>;

private:
  E @\exposidnc{unex}@; // exposition only
};
\end{codeblock}

\begin{itemdescr}
\pnum
\mandates
A program that instantiates the definition of \tcode{unexpected} for a non-object type other than an lvalue reference type, an array type, a specialization of \tcode{unexpected}, or a cv-qualified type is ill-formed.

\pnum
\remarks
Subclause \iref{expected.unexpected} describes the class template \tcode{unexpected} that represents unexpected objects stored in \tcode{expected} objects.
\end{itemdescr}

\rSec3[expected.un.cons]{Constructors}

\indexlibraryctor{unexpected}%
\begin{itemdecl}
template<class Err = E>
  requires(!is_same_v<remove_cvref_t<Err>, unexpected> &&
           !is_same_v<remove_cvref_t<Err>, in_place_t> && is_constructible_v<E, Err>)
constexpr explicit unexpected(Err&& e) noexcept(is_nothrow_constructible_v<E, Err>);
\end{itemdecl}

\begin{itemdescr}
\pnum
\constraints
\begin{itemize}
\item \tcode{is_same_v<remove_cvref_t<Err>, unexpected>} is \tcode{false}; and
\item \tcode{is_same_v<remove_cvref_t<Err>, in_place_t>} is \tcode{false}; and
\item \tcode{is_constructible_v<E, Err>} is \tcode{true}.
\end{itemize}

\pnum
\effects
Direct-non-list-initializes \tcode{unex} with \tcode{std::forward<Err>(e)}.

\pnum
\throws
Any exception thrown by the initialization of \tcode{unex}.
\end{itemdescr}

\indexlibraryctor{unexpected}%
\begin{itemdecl}
template<class... Args>
  requires is_constructible_v<E, Args...>
constexpr explicit unexpected(in_place_t, Args&&... args) noexcept(
    is_nothrow_constructible_v<E, Args...>);
\end{itemdecl}

\begin{itemdescr}
\pnum
\constraints
\tcode{is_constructible_v<E, Args...>} is \tcode{true}.

\pnum
\effects
Direct-non-list-initializes \tcode{unex} with \tcode{std::forward<Args>(args)...}.

\pnum
\throws
Any exception thrown by the initialization of \tcode{unex}.
\end{itemdescr}

\indexlibraryctor{unexpected}%
\begin{itemdecl}
template<class U, class... Args>
  requires is_constructible_v<E, initializer_list<U>&, Args...>
constexpr explicit unexpected(
    in_place_t, initializer_list<U> il,
    Args&&... args) noexcept(is_nothrow_constructible_v<E, initializer_list<U>&,
                                                        Args...>);
\end{itemdecl}

\begin{itemdescr}
\pnum
\constraints
\tcode{is_constructible_v<E, initializer_list<U>&, Args...>} is \tcode{true}.

\pnum
\effects
Direct-non-list-initializes \tcode{unex} with \tcode{il, std::forward<Args>(args)...}.

\pnum
\throws
Any exception thrown by the initialization of \tcode{unex}.
\end{itemdescr}

\rSec3[expected.un.obs]{Observers}

\indexlibrarymember{error}{unexpected}%
\begin{itemdecl}
constexpr const E& error() const& noexcept;
\end{itemdecl}

\begin{itemdescr}
\pnum
\returns
\tcode{unex}.
\end{itemdescr}

\indexlibrarymember{error}{unexpected}%
\begin{itemdecl}
constexpr E& error() & noexcept;
\end{itemdecl}

\indexlibrarymember{error}{unexpected}%
\begin{itemdecl}
constexpr const E&& error() const&& noexcept;
\end{itemdecl}

\begin{itemdescr}
\pnum
\returns
\tcode{std::move(unex)}.
\end{itemdescr}

\indexlibrarymember{error}{unexpected}%
\begin{itemdecl}
constexpr E&& error() && noexcept;
\end{itemdecl}

\rSec3[expected.un.swap]{Swap}

\indexlibrarymember{swap}{unexpected}%
\begin{itemdecl}
constexpr void swap(unexpected& other) noexcept(is_nothrow_swappable_v<E>);
\end{itemdecl}

\begin{itemdescr}
\pnum
\mandates
\tcode{is_swappable_v<E>} is \tcode{true}.

\pnum
\effects
Equivalent to: \tcode{using std::swap; swap(unex, other.unex);}
\end{itemdescr}

\indexlibraryglobal{swap}%
\begin{itemdecl}
friend constexpr void swap(unexpected& x, unexpected& y) noexcept(noexcept(x.swap(y)));
\end{itemdecl}

\begin{itemdescr}
\pnum
\constraints
\tcode{is_swappable_v<E>} is \tcode{true}.

\pnum
\effects
Equivalent to \tcode{x.swap(y)}.
\end{itemdescr}

\rSec3[expected.un.eq]{Equality operator}

\indexlibraryglobal{operator==}%
\begin{itemdecl}
template<class E2>
friend constexpr bool operator==(const unexpected& x, const unexpected<E2>& y);
\end{itemdecl}

\begin{itemdescr}
\pnum
\mandates
The expression \tcode{x.error() == y.error()} is well-formed and its result is convertible to \tcode{bool}.

\pnum
\returns
\tcode{x.error() == y.error()}.
\end{itemdescr}

\rSec3[expected.un.ref]{Partial specialization unexpected<E&>}

\begin{codeblock}
template<class E>
class unexpected<E&> {
public:
  constexpr unexpected(const unexpected&) = default;
  constexpr unexpected(unexpected&&) = default;
  template<class G = E>
    constexpr explicit unexpected(G&&) noexcept;
  template<class G = E>
    constexpr explicit unexpected(in_place_t, G&&) noexcept;

  constexpr unexpected& operator=(const unexpected&) = default;
  constexpr unexpected& operator=(unexpected&&) = default;

  constexpr E& error() const noexcept;

  constexpr void swap(unexpected& other) noexcept;

  template<class E2>
    friend constexpr bool operator==(const unexpected&, const unexpected<E2>&);

  friend constexpr void swap(unexpected& x, unexpected& y) noexcept;

private:
  E* unex;              // exposition only
};
\end{codeblock}

\begin{itemdescr}
\pnum
\mandates
A program that instantiates the definition of \tcode{unexpected<E&>} for an array type or a specialization of \tcode{unexpected} is ill-formed.

\pnum
\remarks
An object of type \tcode{unexpected<E&>} holds a pointer to an object of type \tcode{E}. The referenced object is not owned by the \tcode{unexpected<E&>} object. Unlike the primary template, \tcode{E} may be a cv-qualified type.
\end{itemdescr}

\indexlibraryctor{unexpected}%
\begin{itemdecl}
template<class G = E>
  requires(!is_same_v<remove_cvref_t<G>, unexpected> &&
           !is_same_v<remove_cvref_t<G>, in_place_t> && is_constructible_v<E&, G&&> &&
           !reference_constructs_from_temporary_v<E&, G>)
constexpr explicit unexpected(G&& e) noexcept;
\end{itemdecl}

\begin{itemdescr}
\pnum
\constraints
\begin{itemize}
\item \tcode{is_same_v<remove_cvref_t<G>, unexpected>} is \tcode{false},
\item \tcode{is_same_v<remove_cvref_t<G>, in_place_t>} is \tcode{false},
\item \tcode{is_constructible_v<E&, G>} is \tcode{true}, and
\item \tcode{reference_constructs_from_temporary_v<E&, G>} is \tcode{false}.
\end{itemize}

\pnum
\effects
Initializes \tcode{unex} with \tcode{addressof(static_cast<E&>(std::forward<G>(e)))}.

\pnum
\remarks
A constructor for which \tcode{reference_constructs_from_temporary_v<E&, G>} is \tcode{true} — one that would bind \tcode{E&} to a temporary — is defined as deleted.
\end{itemdescr}

\indexlibraryctor{unexpected}%
\begin{itemdecl}
template<class G = E>
  requires(is_constructible_v<E&, G&&> && !reference_constructs_from_temporary_v<E&, G>)
constexpr explicit unexpected(in_place_t, G&& e) noexcept;
\end{itemdecl}

\begin{itemdescr}
\pnum
\constraints
\tcode{is_constructible_v<E&, G>} is \tcode{true} and \tcode{reference_constructs_from_temporary_v<E&, G>} is \tcode{false}.

\pnum
\effects
Initializes \tcode{unex} with \tcode{addressof(static_cast<E&>(std::forward<G>(e)))}.
\end{itemdescr}

\indexlibrarymember{error}{unexpected}%
\begin{itemdecl}
constexpr E& error() const noexcept;
\end{itemdecl}

\begin{itemdescr}
\pnum
\returns
\tcode{*unex}.

\pnum
\remarks
The reference returned is not \tcode{const}-qualified even when \tcode{*this} is \tcode{const}; the constness of the \tcode{unexpected<E&>} object does not propagate to the referenced object.
\end{itemdescr}

\indexlibrarymember{swap}{unexpected}%
\begin{itemdecl}
constexpr void swap(unexpected& other) noexcept;
\end{itemdecl}

\begin{itemdescr}
\pnum
\effects
Exchanges \tcode{unex} and \tcode{other.unex}.
\end{itemdescr}

\indexlibraryglobal{operator==}%
\begin{itemdecl}
template<class E2>
friend constexpr bool operator==(const unexpected& x, const unexpected<E2>& y);
\end{itemdecl}

\begin{itemdescr}
\pnum
\mandates
The expression \tcode{x.error() == y.error()} is well-formed and its result is convertible to \tcode{bool}.

\pnum
\returns
\tcode{x.error() == y.error()}.
\end{itemdescr}

\indexlibraryglobal{swap}%
\begin{itemdecl}
friend constexpr void swap(unexpected& x, unexpected& y) noexcept;
\end{itemdecl}

\begin{itemdescr}
\pnum
\effects
Equivalent to \tcode{x.swap(y)}.
\end{itemdescr}

\rSec2[expected.bad]{Class template bad_expected_access}

\indexlibrarymember{what}{bad_expected_access}%
\begin{itemdecl}
const char* what() const noexcept override;
\end{itemdecl}

\begin{itemdescr}
\pnum
\returns
An implementation-defined ntbs, which during constant evaluation is encoded with the ordinary literal encoding (\iref{lex.ccon}).
\end{itemdescr}

\begin{codeblock}
template<class E>
class @\libglobal{bad_expected_access}@ : public bad_expected_access<void> {
public:
  BEMAN_EXPECTED_CONSTEXPR_EXCEPTION explicit bad_expected_access(E e);
  BEMAN_EXPECTED_CONSTEXPR_EXCEPTION const char* what() const noexcept override;
  constexpr E& error() & noexcept;
  constexpr const E& error() const& noexcept;
  constexpr E&& error() && noexcept;
  constexpr const E&& error() const&& noexcept;

private:
  E @\exposidnc{unex}@; // exposition only
};
\end{codeblock}

\begin{itemdescr}
\pnum
\remarks
The class template \tcode{bad_expected_access} defines the type of objects thrown as exceptions to report the situation where an attempt is made to access the value of an \tcode{expected<T, E>} object for which \tcode{has_value()} is \tcode{false}.
\end{itemdescr}

\indexlibraryctor{bad_expected_access}%
\begin{itemdecl}
explicit bad_expected_access(E e);
\end{itemdecl}

\begin{itemdescr}
\pnum
\effects
Initializes \exposid{unex} with \tcode{std::move(e)}.
\end{itemdescr}

\indexlibrarymember{what}{bad_expected_access}%
\begin{itemdecl}
const char* what() const noexcept override;
\end{itemdecl}

\begin{itemdescr}
\pnum
\returns
An implementation-defined ntbs, which during constant evaluation is encoded with the ordinary literal encoding (\iref{lex.ccon}).
\end{itemdescr}

\indexlibrarymember{error}{bad_expected_access}%
\begin{itemdecl}
constexpr E& error() & noexcept;
\end{itemdecl}

\begin{itemdescr}
\pnum
\returns
\exposid{unex}.
\end{itemdescr}

\indexlibrarymember{error}{bad_expected_access}%
\begin{itemdecl}
constexpr const E& error() const& noexcept;
\end{itemdecl}

\begin{itemdescr}
\pnum
\returns
\exposid{unex}.
\end{itemdescr}

\indexlibrarymember{error}{bad_expected_access}%
\begin{itemdecl}
constexpr E&& error() && noexcept;
\end{itemdecl}

\begin{itemdescr}
\pnum
\returns
\tcode{std::move(unex)}.
\end{itemdescr}

\indexlibrarymember{error}{bad_expected_access}%
\begin{itemdecl}
constexpr const E&& error() const&& noexcept;
\end{itemdecl}

\begin{itemdescr}
\pnum
\returns
\tcode{std::move(unex)}.
\end{itemdescr}

\rSec2[expected.bad.void]{Class template specialization bad_expected_access<void>}

\begin{codeblock}
template<>
class @\libglobal{bad_expected_access}@<void> : public exception {
protected:
  BEMAN_EXPECTED_CONSTEXPR_EXCEPTION bad_expected_access() noexcept = default;
  BEMAN_EXPECTED_CONSTEXPR_EXCEPTION bad_expected_access(
      const bad_expected_access&) noexcept = default;
  BEMAN_EXPECTED_CONSTEXPR_EXCEPTION bad_expected_access(
      bad_expected_access&&) noexcept = default;
  BEMAN_EXPECTED_CONSTEXPR_EXCEPTION bad_expected_access& operator=(
      const bad_expected_access&) noexcept = default;
  BEMAN_EXPECTED_CONSTEXPR_EXCEPTION bad_expected_access& operator=(
      bad_expected_access&&) noexcept = default;
  BEMAN_EXPECTED_CONSTEXPR_EXCEPTION ~bad_expected_access() = default;

public:
  BEMAN_EXPECTED_CONSTEXPR_EXCEPTION const char* what() const noexcept override;
};
\end{codeblock}

\rSec2[expected.expected]{Class template expected}

\rSec3[expected.object.general]{General}

\begin{codeblock}
template<class T, class E>
class @\libglobal{expected}@ {
private:
  using @\exposidnc{error-value-type}@ = remove_cv_t<remove_reference_t<E>>; // exposition only

public:
  using @\libmember{value_type}{expected}@ = T;
  using @\libmember{error_type}{expected}@ = E;
  using @\libmember{unexpected_type}{expected}@ = unexpected<E>;

  template<class U> using rebind = expected<U, error_type>;

  // -------------------------------------------------------------------------
  // [expected.object.cons] Constructors
  // -------------------------------------------------------------------------

  // Default constructor: value-initializes T
  constexpr expected() noexcept(is_nothrow_default_constructible_v<T>)
    requires is_default_constructible_v<T>;

  // Copy constructor (trivial path). Unconstrained on purpose: this is the
  // sole declaration when T or E is not copy constructible at all (it is
  // then implicitly defined as deleted), and the more-constrained
  // non-trivial-path overload below is selected over it by constraint
  // subsumption whenever it is viable.
  constexpr expected(const expected&) = default;

  // Copy constructor (non-trivial path)
  constexpr expected(const expected& rhs) noexcept(is_nothrow_copy_constructible_v<T> &&
                                                   is_nothrow_copy_constructible_v<E>)
    requires(is_copy_constructible_v<T> && is_copy_constructible_v<E> &&
             !(is_trivially_copy_constructible_v<T> &&
               is_trivially_copy_constructible_v<E>));

  // Move constructor (trivial path). Unconstrained; see the copy
  // constructor above for why. No explicit noexcept: let the compiler
  // deduce it, so a non-movable-at-all E deletes rather than mismatches.
  constexpr expected(expected&&) = default;

  // Move constructor (non-trivial path)
  constexpr expected(expected&& rhs) noexcept(is_nothrow_move_constructible_v<T> &&
                                              is_nothrow_move_constructible_v<E>)
    requires(is_move_constructible_v<T> && is_move_constructible_v<E> &&
             !(is_trivially_move_constructible_v<T> &&
               is_trivially_move_constructible_v<E>));

  // Converting copy constructor from expected<U, G> — value-E path
  template<class U, class G>
    requires(!is_reference_v<E> && is_constructible_v<T, const U&> &&
             is_constructible_v<E, const G&> &&
             (is_same_v<bool, remove_cv_t<T>> ||
              !@\exposidnc{converts-from-any-cvref}@<T, expected<U, G>>) &&
             !is_constructible_v<unexpected<E>, expected<U, G>&> &&
             !is_constructible_v<unexpected<E>, expected<U, G>&&> &&
             !is_constructible_v<unexpected<E>, const expected<U, G>&> &&
             !is_constructible_v<unexpected<E>, const expected<U, G>&&>)
  constexpr explicit(!is_convertible_v<const U&, T> || !is_convertible_v<const G&, E>)
      expected(const expected<U, G>& rhs);

  // Converting move constructor from expected<U, G> — value-E path
  template<class U, class G>
    requires(!is_reference_v<E> && is_constructible_v<T, U> &&
             is_constructible_v<E, G> &&
             (is_same_v<bool, remove_cv_t<T>> ||
              !@\exposidnc{converts-from-any-cvref}@<T, expected<U, G>>) &&
             !is_constructible_v<unexpected<E>, expected<U, G>&> &&
             !is_constructible_v<unexpected<E>, expected<U, G>&&> &&
             !is_constructible_v<unexpected<E>, const expected<U, G>&> &&
             !is_constructible_v<unexpected<E>, const expected<U, G>&&>)
  constexpr explicit(!is_convertible_v<U, T> || !is_convertible_v<G, E>)
      expected(expected<U, G>&& rhs);

  // Converting constructor from expected<U, G> — reference-E path: only accepts sources
  // whose error type G is itself a reference convertible to E.
  template<class U, class G>
    requires(is_reference_v<E> && is_reference_v<G> &&
             is_constructible_v<T, const U&> && is_convertible_v<G, E>)
  constexpr explicit(!is_convertible_v<const U&, T> || !is_convertible_v<G, E>)
      expected(const expected<U, G>& rhs);

  template<class U, class G>
    requires(is_reference_v<E> && is_reference_v<G> && is_constructible_v<T, U&&> &&
             is_convertible_v<G, E>)
  constexpr explicit(!is_convertible_v<U&&, T> || !is_convertible_v<G, E>)
      expected(expected<U, G>&& rhs);

  // Constructor from value U&&
  template<class U = remove_cv_t<T>>
    requires(!is_same_v<remove_cvref_t<U>, in_place_t> &&
             !is_same_v<remove_cvref_t<U>, unexpect_t> &&
             !is_same_v<remove_cvref_t<U>, expected> && is_constructible_v<T, U> &&
             !@\exposidnc{is-unexpected-specialization}@<remove_cvref_t<U>>::value &&
             (!is_same_v<bool, remove_cv_t<T>> ||
              !@\exposidnc{is-expected-specialization}@<remove_cvref_t<U>>::value))
  constexpr explicit(!is_convertible_v<U, T>) expected(U&& v);

  // Constructor from unexpected<G> const& / && — value-E path
  template<class G>
    requires(!is_reference_v<E> && is_constructible_v<E, const G&>)
  constexpr explicit(!is_convertible_v<const G&, E>) expected(const unexpected<G>& e);

  template<class G>
    requires(!is_reference_v<E> && is_constructible_v<E, G>)
  constexpr explicit(!is_convertible_v<G, E>) expected(unexpected<G>&& e);

  // Constructor from unexpected<G> — reference-E path. Allowed only when G is itself a
  // reference, i.e. the source unexpected holds a reference to an external object, so
  // binding E& to e.error() cannot dangle regardless of the source's value category. No
  // const_cast is needed.
  template<class G>
    requires(is_reference_v<E> && is_reference_v<G> && is_constructible_v<E, G> &&
             !reference_constructs_from_temporary_v<E, G>)
  constexpr explicit(!is_convertible_v<G, E>) expected(const unexpected<G>& e) noexcept;

  template<class G>
    requires(is_reference_v<E> && is_reference_v<G> && is_constructible_v<E, G> &&
             !reference_constructs_from_temporary_v<E, G>)
  constexpr explicit(!is_convertible_v<G, E>) expected(unexpected<G>&& e) noexcept;

  // Deleted for reference E with value G: the referent lives inside the unexpected<G>
  // object, so binding E& to it would dangle once a temporary source is destroyed. Use
  // (unexpect, lvalue), or an unexpected<E&> holding an external object, instead.
  template<class G>
    requires(is_reference_v<E> && !is_reference_v<G>)
  constexpr expected(const unexpected<G>&) = BEMAN_EXPECTED_DELETE_MSG(
      "expected<T,E&>: cannot construct from unexpected<value>; the value would dangle "
      "— use unexpected<E&>");

  template<class G>
    requires(is_reference_v<E> && !is_reference_v<G>)
  constexpr expected(unexpected<G>&&) = BEMAN_EXPECTED_DELETE_MSG(
      "expected<T,E&>: cannot construct from unexpected<value>; the value would dangle "
      "— use unexpected<E&>");

  // In-place constructor for value
  template<class... Args>
    requires is_constructible_v<T, Args...>
  constexpr explicit expected(in_place_t, Args&&... args);

  // In-place constructor for value with initializer_list
  template<class U, class... Args>
    requires is_constructible_v<T, initializer_list<U>&, Args...>
  constexpr explicit expected(in_place_t, initializer_list<U> il, Args&&... args);

  // In-place constructor for error
  template<class... Args>
    requires(is_constructible_v<E, Args...> && !@\exposidnc{unexpect-dangles-v}@<E, Args...>)
  constexpr explicit expected(unexpect_t, Args&&... args);

  // Deleted: single argument would bind E& to a temporary — dangling prevention
  template<class... Args>
    requires(@\exposidnc{unexpect-dangles-v}@<E, Args...>)
  constexpr expected(unexpect_t, Args&&...) =
      BEMAN_EXPECTED_DELETE_MSG("expected<T,E&>: unexpect argument would bind a "
                                "temporary that dangles; pass an lvalue reference");

  // Deleted catch-all: reference E, argument neither constructible nor a dangling case
  // (e.g. binding a non-const E& from a const lvalue).
  template<class... Args>
    requires(is_reference_v<E> && !is_constructible_v<E, Args...> &&
             !@\exposidnc{unexpect-dangles-v}@<E, Args...>)
  constexpr expected(unexpect_t, Args&&...) = BEMAN_EXPECTED_DELETE_MSG(
      "expected<T,E&>: no viable conversion from the given argument(s) to E&");

  // In-place constructor for error with initializer_list
  template<class U, class... Args>
    requires(!is_reference_v<E> && is_constructible_v<E, initializer_list<U>&, Args...>)
  constexpr explicit expected(unexpect_t, initializer_list<U> il, Args&&... args);

  template<class U, class... Args>
    requires is_reference_v<E>
  constexpr expected(unexpect_t, initializer_list<U>, Args&&...) =
      BEMAN_EXPECTED_DELETE_MSG("expected<T,E&>: initializer-list error construction "
                                "cannot bind a reference; pass an lvalue reference");

  // -------------------------------------------------------------------------
  // [expected.object.dtor] Destructor
  // -------------------------------------------------------------------------

  constexpr ~expected()
    requires(is_trivially_destructible_v<T> && is_trivially_destructible_v<E>)
  = default;

  constexpr ~expected()
    requires(!(is_trivially_destructible_v<T> && is_trivially_destructible_v<E>));

  // -------------------------------------------------------------------------
  // [expected.object.assign] Assignment
  // -------------------------------------------------------------------------

  // Copy assignment (trivial path)
  constexpr expected& operator=(const expected&)
    requires(is_trivially_copy_constructible_v<T> &&
             is_trivially_copy_assignable_v<T> && is_trivially_destructible_v<T> &&
             is_trivially_copy_constructible_v<E> &&
             is_trivially_copy_assignable_v<E> && is_trivially_destructible_v<E>)
  = default;

  // Copy assignment (non-trivial path)
  constexpr expected& operator=(const expected& rhs) noexcept(
      is_nothrow_copy_constructible_v<T> && is_nothrow_copy_assignable_v<T> &&
      is_nothrow_copy_constructible_v<E> && is_nothrow_copy_assignable_v<E>)
    requires(is_copy_constructible_v<T> && is_copy_assignable_v<T> &&
             (is_reference_v<E> ||
              (is_copy_constructible_v<E> && is_copy_assignable_v<E>)) &&
             (is_nothrow_move_constructible_v<T> ||
              is_nothrow_move_constructible_v<E>) &&
             !(is_trivially_copy_constructible_v<T> &&
               is_trivially_copy_assignable_v<T> && is_trivially_destructible_v<T> &&
               is_trivially_copy_constructible_v<E> &&
               is_trivially_copy_assignable_v<E> && is_trivially_destructible_v<E>));

  // Move assignment (trivial path)
  constexpr expected& operator=(expected&&) noexcept
    requires(is_trivially_move_constructible_v<T> &&
             is_trivially_move_assignable_v<T> && is_trivially_destructible_v<T> &&
             is_trivially_move_constructible_v<E> &&
             is_trivially_move_assignable_v<E> && is_trivially_destructible_v<E>)
  = default;

  // Move assignment (non-trivial path)
  constexpr expected& operator=(expected&& rhs) noexcept(
      is_nothrow_move_constructible_v<T> && is_nothrow_move_assignable_v<T> &&
      is_nothrow_move_constructible_v<E> && is_nothrow_move_assignable_v<E>)
    requires(is_move_constructible_v<T> && is_move_assignable_v<T> &&
             (is_reference_v<E> ||
              (is_move_constructible_v<E> && is_move_assignable_v<E>)) &&
             (is_nothrow_move_constructible_v<T> ||
              is_nothrow_move_constructible_v<E>) &&
             !(is_trivially_move_constructible_v<T> &&
               is_trivially_move_assignable_v<T> && is_trivially_destructible_v<T> &&
               is_trivially_move_constructible_v<E> &&
               is_trivially_move_assignable_v<E> && is_trivially_destructible_v<E>));

  // Assignment from value U&&
  template<class U = remove_cv_t<T>>
    requires(!is_same_v<expected, remove_cvref_t<U>> &&
             !@\exposidnc{is-unexpected-specialization}@<remove_cvref_t<U>>::value &&
             is_constructible_v<T, U> && is_assignable_v<T&, U> &&
             (is_nothrow_constructible_v<T, U> || is_nothrow_move_constructible_v<T> ||
              is_nothrow_move_constructible_v<E>))
  constexpr expected& operator=(U&& v);

  // Assignment from unexpected<G> — value-E path
  template<class G>
    requires(!is_reference_v<E> && is_constructible_v<E, const G&> &&
             is_assignable_v<E&, const G&> &&
             (is_nothrow_constructible_v<E, const G&> ||
              is_nothrow_move_constructible_v<T> || is_nothrow_move_constructible_v<E>))
  constexpr expected& operator=(const unexpected<G>& e);

  template<class G>
    requires(!is_reference_v<E> && is_constructible_v<E, G> && is_assignable_v<E&, G> &&
             (is_nothrow_constructible_v<E, G> || is_nothrow_move_constructible_v<T> ||
              is_nothrow_move_constructible_v<E>))
  constexpr expected& operator=(unexpected<G>&& e);

  // Rebinding assignment for reference E from reference G — binds unex_ to the external
  // referent (never dangles; pointer store is noexcept). Mirrors the reference-E
  // constructor.
  template<class G>
    requires(is_reference_v<E> && is_reference_v<G> && is_constructible_v<E, G> &&
             !reference_constructs_from_temporary_v<E, G>)
  constexpr expected& operator=(const unexpected<G>& e);

  template<class G>
    requires(is_reference_v<E> && is_reference_v<G> && is_constructible_v<E, G> &&
             !reference_constructs_from_temporary_v<E, G>)
  constexpr expected& operator=(unexpected<G>&& e);

  // Deleted for reference E with value G: would rebind E& to unexpected<G>'s temporary
  // storage.
  template<class G>
    requires(is_reference_v<E> && !is_reference_v<G>)
  constexpr expected& operator=(const unexpected<G>&) =
      BEMAN_EXPECTED_DELETE_MSG("expected<T,E&>: cannot assign from unexpected<value>; "
                                "the value would dangle — use unexpected<E&>");

  template<class G>
    requires(is_reference_v<E> && !is_reference_v<G>)
  constexpr expected& operator=(unexpected<G>&&) =
      BEMAN_EXPECTED_DELETE_MSG("expected<T,E&>: cannot assign from unexpected<value>; "
                                "the value would dangle — use unexpected<E&>");

  // Emplace: destroy current value/error, construct value in-place
  template<class... Args>
    requires is_nothrow_constructible_v<T, Args...>
  constexpr T& emplace(Args&&... args) noexcept;

  template<class U, class... Args>
    requires is_nothrow_constructible_v<T, initializer_list<U>&, Args...>
  constexpr T& emplace(initializer_list<U> il, Args&&... args) noexcept;

  // -------------------------------------------------------------------------
  // [expected.object.swap] Swap
  // -------------------------------------------------------------------------

  constexpr void swap(expected& rhs) noexcept(is_nothrow_move_constructible_v<T> &&
                                              is_nothrow_swappable_v<T> &&
                                              is_nothrow_move_constructible_v<E> &&
                                              (is_reference_v<E> ||
                                               is_nothrow_swappable_v<E>))
    requires(is_swappable_v<T> && (is_reference_v<E> || is_swappable_v<E>) &&
             is_move_constructible_v<T> && is_move_constructible_v<E> &&
             (is_nothrow_move_constructible_v<T> || is_nothrow_move_constructible_v<E>))
  ;

  friend constexpr void swap(expected& x, expected& y) noexcept(noexcept(x.swap(y)))
    requires(is_swappable_v<T> && (is_reference_v<E> || is_swappable_v<E>) &&
             is_move_constructible_v<T> && is_move_constructible_v<E> &&
             (is_nothrow_move_constructible_v<T> || is_nothrow_move_constructible_v<E>))
  ;

  // -------------------------------------------------------------------------
  // [expected.object.obs] Observers
  // -------------------------------------------------------------------------

  constexpr const T* operator->() const noexcept;
  constexpr T* operator->() noexcept;

  constexpr const T& operator*() const& noexcept;
  constexpr T& operator*() & noexcept;
  constexpr const T&& operator*() const&& noexcept;
  constexpr T&& operator*() && noexcept;

  constexpr explicit operator bool() const noexcept;
  constexpr bool has_value() const noexcept;

  constexpr const T& value() const&;
  constexpr T& value() &;
  constexpr const T&& value() const&&;
  constexpr T&& value() &&;

  constexpr const E& error() const& noexcept;
  constexpr E& error() & noexcept;
  constexpr const E&& error() const&& noexcept;
  constexpr E&& error() && noexcept;

  template<class U = remove_cv_t<T>> constexpr T value_or(U&& def) const&;

  template<class U = remove_cv_t<T>> constexpr T value_or(U&& def) &&;

  template<class G = @\exposidnc{error-value-type}@>
    requires(is_copy_constructible_v<@\exposidnc{error-value-type}@> &&
             is_convertible_v<G, @\exposidnc{error-value-type}@>)
  constexpr @\exposidnc{error-value-type}@ error_or(G&& def) const&;

  template<class G = @\exposidnc{error-value-type}@>
    requires(is_move_constructible_v<@\exposidnc{error-value-type}@> &&
             is_convertible_v<G, @\exposidnc{error-value-type}@>)
  constexpr @\exposidnc{error-value-type}@ error_or(G&& def) &&;

  // -------------------------------------------------------------------------
  // [expected.object.monadic] Monadic operations
  // -------------------------------------------------------------------------

  template<class F>
    requires is_constructible_v<E, E&>
  constexpr auto and_then(F&& f) &;
  template<class F>
    requires is_constructible_v<E, E&&>
  constexpr auto and_then(F&& f) &&;
  template<class F>
    requires is_constructible_v<E, const E&>
  constexpr auto and_then(F&& f) const&;
  template<class F>
    requires is_constructible_v<E, const E&&>
  constexpr auto and_then(F&& f) const&&;

  template<class F>
    requires is_constructible_v<T, T&>
  constexpr auto or_else(F&& f) &;
  template<class F>
    requires is_constructible_v<T, T&&>
  constexpr auto or_else(F&& f) &&;
  template<class F>
    requires is_constructible_v<T, const T&>
  constexpr auto or_else(F&& f) const&;
  template<class F>
    requires is_constructible_v<T, const T&&>
  constexpr auto or_else(F&& f) const&&;

  template<class F>
    requires is_constructible_v<E, E&>
  constexpr auto transform(F&& f) &;
  template<class F>
    requires is_constructible_v<E, E&&>
  constexpr auto transform(F&& f) &&;
  template<class F>
    requires is_constructible_v<E, const E&>
  constexpr auto transform(F&& f) const&;
  template<class F>
    requires is_constructible_v<E, const E&&>
  constexpr auto transform(F&& f) const&&;

  template<class F>
    requires is_constructible_v<T, T&>
  constexpr auto transform_error(F&& f) &;
  template<class F>
    requires is_constructible_v<T, T&&>
  constexpr auto transform_error(F&& f) &&;
  template<class F>
    requires is_constructible_v<T, const T&>
  constexpr auto transform_error(F&& f) const&;
  template<class F>
    requires is_constructible_v<T, const T&&>
  constexpr auto transform_error(F&& f) const&&;

  // -------------------------------------------------------------------------
  // [expected.object.eq] Equality operators (hidden friends)
  // -------------------------------------------------------------------------

  template<class T2, class E2>
    requires(!is_void_v<T2>)
  friend constexpr bool operator==(const expected& x, const expected<T2, E2>& y);

  template<class T2>
    requires(!@\exposidnc{is-expected-specialization}@<T2>::value)
  friend constexpr bool operator==(const expected& x, const T2& val);

  template<class E2>
  friend constexpr bool operator==(const expected& x, const unexpected<E2>& e);

private:
  bool @\exposidnc{has-val}@; // exposition only
  union {
    T @\exposidnc{val}@;              // exposition only
    unexpected<E> @\exposidnc{unex}@; // exposition only
  };
};
\end{codeblock}

\begin{itemdescr}
\pnum
\mandates
A program that instantiates the definition of \tcode{expected<T, E>} with a \tcode{T} that is not a valid value type for \tcode{expected} (that is, \tcode{remove_cv_t<T>} is \tcode{void}, or a complete non-array object type other than \tcode{in_place_t}, \tcode{unexpect_t}, or a specialization of \tcode{unexpected}) is ill-formed. A program that instantiates the definition of \tcode{expected<T, E>} with an \tcode{E} that is not a valid template argument for \tcode{unexpected} is ill-formed.

\pnum
\remarks
Any object of type \tcode{expected<T, E>} either contains a value of type \tcode{T} or a value of type \tcode{E} nested within it. Member \tcode{has_val} indicates whether the \tcode{expected<T, E>} object contains an object of type \tcode{T}. When \tcode{has_value()} is \tcode{false}, the error is \tcode{unex.error()}. The error is held as an \tcode{unexpected<E>}, and not as an \tcode{E}, so that \tcode{E} may be an lvalue reference type: an \tcode{E&} cannot be a union member, whereas \tcode{unexpected<E&>} holds a pointer to an external object.
\end{itemdescr}

\rSec3[expected.object.cons]{Constructors}

\indexlibraryctor{expected}%
\begin{itemdecl}
constexpr expected() noexcept(is_nothrow_default_constructible_v<T>);
\end{itemdecl}

\begin{itemdescr}
\pnum
\constraints
\tcode{is_default_constructible_v<T>} is \tcode{true}.

\pnum
\effects
Value-initializes \tcode{val}.

\pnum
\ensures
\tcode{has_value()} is \tcode{true}.

\pnum
\throws
Any exception thrown by the initialization of \tcode{val}.
\end{itemdescr}

\indexlibraryctor{expected}%
\begin{itemdecl}
constexpr expected(const expected&) = default;
\end{itemdecl}

\begin{itemdescr}
\pnum
\effects
Direct-non-list-initializes \tcode{val} or \tcode{unex} (matching \tcode{rhs}'s active member) by trivial copy construction.

\pnum
\ensures
\tcode{rhs.has_value() == this->has_value()}.

\pnum
\remarks
This constructor is trivial.
\end{itemdescr}

\indexlibraryctor{expected}%
\begin{itemdecl}
constexpr expected(const expected& rhs) noexcept(is_nothrow_copy_constructible_v<T> &&
                                                 is_nothrow_copy_constructible_v<E>);
\end{itemdecl}

\begin{itemdescr}
\pnum
\constraints
\tcode{is_copy_constructible_v<T>} is \tcode{true}, \tcode{is_copy_constructible_v<E>} is \tcode{true}, and \tcode{(is_trivially_copy_constructible_v<T> && is_trivially_copy_constructible_v<E>)} is \tcode{false}.

\pnum
\effects
If \tcode{rhs.has_value()} is \tcode{true}, direct-non-list-initializes \tcode{val} with \tcode{*rhs}. Otherwise, direct-non-list-initializes \tcode{unex} with \tcode{rhs.error()}.

\pnum
\ensures
\tcode{rhs.has_value() == this->has_value()}.

\pnum
\throws
Any exception thrown by the initialization of \tcode{val} or \tcode{unex}.

\pnum
\remarks
This constructor is defined as deleted unless \tcode{is_copy_constructible_v<T>} is \tcode{true} and \tcode{is_copy_constructible_v<E>} is \tcode{true} or \tcode{is_reference_v<E>} is \tcode{true}. This constructor is trivial if \tcode{is_trivially_copy_constructible_v<T>} is \tcode{true} and \tcode{is_trivially_copy_constructible_v<E>} is \tcode{true}.
\end{itemdescr}

\indexlibraryctor{expected}%
\begin{itemdecl}
constexpr expected(expected&&) = default;
\end{itemdecl}

\begin{itemdescr}
\pnum
\effects
Direct-non-list-initializes \tcode{val} or \tcode{unex} (matching \tcode{rhs}'s active member) by trivial move construction.

\pnum
\ensures
\tcode{rhs.has_value() == this->has_value()}.

\pnum
\remarks
This constructor is trivial.
\end{itemdescr}

\indexlibraryctor{expected}%
\begin{itemdecl}
constexpr expected(expected&& rhs) noexcept(is_nothrow_move_constructible_v<T> &&
                                            is_nothrow_move_constructible_v<E>);
\end{itemdecl}

\begin{itemdescr}
\pnum
\constraints
\tcode{is_move_constructible_v<T>} is \tcode{true}, \tcode{is_move_constructible_v<E>} is \tcode{true}, and \tcode{(is_trivially_move_constructible_v<T> && is_trivially_move_constructible_v<E>)} is \tcode{false}.

\pnum
\effects
If \tcode{rhs.has_value()} is \tcode{true}, direct-non-list-initializes \tcode{val} with \tcode{std::move(*rhs)}. Otherwise, direct-non-list-initializes \tcode{unex} with \tcode{std::move(rhs.error())}.

\pnum
\ensures
\tcode{rhs.has_value()} is unchanged; \tcode{rhs.has_value() == this->has_value()} is \tcode{true}.

\pnum
\throws
Any exception thrown by the initialization of \tcode{val} or \tcode{unex}.

\pnum
\remarks
This constructor is trivial if \tcode{is_trivially_move_constructible_v<T>} is \tcode{true} and \tcode{is_trivially_move_constructible_v<E>} is \tcode{true}.
\end{itemdescr}

\indexlibraryctor{expected}%
\begin{itemdecl}
template<class U, class G>
  requires(!is_reference_v<E> && is_constructible_v<T, const U&> &&
           is_constructible_v<E, const G&> &&
           (is_same_v<bool, remove_cv_t<T>> ||
            !@\exposidnc{converts-from-any-cvref}@<T, expected<U, G>>) &&
           !is_constructible_v<unexpected<E>, expected<U, G>&> &&
           !is_constructible_v<unexpected<E>, expected<U, G>&&> &&
           !is_constructible_v<unexpected<E>, const expected<U, G>&> &&
           !is_constructible_v<unexpected<E>, const expected<U, G>&&>)
constexpr explicit(!is_convertible_v<const U&, T> || !is_convertible_v<const G&, E>)
    expected(const expected<U, G>& rhs);
template<class U, class G>
  requires(!is_reference_v<E> && is_constructible_v<T, U> && is_constructible_v<E, G> &&
           (is_same_v<bool, remove_cv_t<T>> ||
            !@\exposidnc{converts-from-any-cvref}@<T, expected<U, G>>) &&
           !is_constructible_v<unexpected<E>, expected<U, G>&> &&
           !is_constructible_v<unexpected<E>, expected<U, G>&&> &&
           !is_constructible_v<unexpected<E>, const expected<U, G>&> &&
           !is_constructible_v<unexpected<E>, const expected<U, G>&&>)
constexpr explicit(!is_convertible_v<U, T> || !is_convertible_v<G, E>)
    expected(expected<U, G>&& rhs);
\end{itemdecl}

\begin{itemdescr}
\pnum
\constraints
\tcode{is_constructible_v<T, const U&>} is \tcode{true}; and \tcode{is_constructible_v<E, const G&>} is \tcode{true}; and if \tcode{T} is not \tcode{bool}, \tcode{converts-from-any-cvref<T, expected<U, G>>} is \tcode{false}; and \tcode{is_constructible_v<unexpected<E>, expected<U, G>&>} is \tcode{false}; and \tcode{is_constructible_v<unexpected<E>, expected<U, G>>} is \tcode{false}; and \tcode{is_constructible_v<unexpected<E>, const expected<U, G>&>} is \tcode{false}; and \tcode{is_constructible_v<unexpected<E>, const expected<U, G>>} is \tcode{false}.

\pnum
\effects
If \tcode{rhs.has_value()}, direct-non-list-initializes \tcode{val} with \tcode{*rhs}. Otherwise, direct-non-list-initializes \tcode{unex} with \tcode{rhs.error()}.

\pnum
\ensures
\tcode{rhs.has_value()} is unchanged; \tcode{rhs.has_value() == this->has_value()} is \tcode{true}.

\pnum
\throws
Any exception thrown by the initialization of \tcode{val} or \tcode{unex}.
\end{itemdescr}

\indexlibraryctor{expected}%
\begin{itemdecl}
template<class U, class G>
  requires(is_reference_v<E> && is_reference_v<G> && is_constructible_v<T, const U&> &&
           is_convertible_v<G, E>)
constexpr explicit(!is_convertible_v<const U&, T> || !is_convertible_v<G, E>)
    expected(const expected<U, G>& rhs);
template<class U, class G>
  requires(is_reference_v<E> && is_reference_v<G> && is_constructible_v<T, U&&> &&
           is_convertible_v<G, E>)
constexpr explicit(!is_convertible_v<U&&, T> || !is_convertible_v<G, E>)
    expected(expected<U, G>&& rhs);
\end{itemdecl}

\begin{itemdescr}
\pnum
\constraints
\tcode{is_reference_v<G>} is \tcode{true} and \tcode{is_convertible_v<G, E>} is \tcode{true}.

\pnum
\effects
If \tcode{rhs.has_value()}, direct-non-list-initializes \tcode{val} with \tcode{*rhs}. Otherwise, direct-non-list-initializes \tcode{unex} with \tcode{rhs.error()}.

\pnum
\ensures
\tcode{rhs.has_value()} is unchanged; \tcode{rhs.has_value() == this->has_value()} is \tcode{true}.

\pnum
\throws
Any exception thrown by the initialization of \tcode{val} or \tcode{unex}.

\pnum
\remarks
Unlike the value-\tcode{E} overload above, this overload participates in overload resolution only when \tcode{E} and \tcode{G} are both reference types, so the referenced error object is never copied.
\end{itemdescr}

\indexlibraryctor{expected}%
\begin{itemdecl}
template<class U = remove_cv_t<T>>
  requires(!is_same_v<remove_cvref_t<U>, in_place_t> &&
           !is_same_v<remove_cvref_t<U>, unexpect_t> &&
           !is_same_v<remove_cvref_t<U>, expected> && is_constructible_v<T, U> &&
           !@\exposidnc{is-unexpected-specialization}@<remove_cvref_t<U>>::value &&
           (!is_same_v<bool, remove_cv_t<T>> ||
            !@\exposidnc{is-expected-specialization}@<remove_cvref_t<U>>::value))
constexpr explicit(!is_convertible_v<U, T>) expected(U&& v);
\end{itemdecl}

\begin{itemdescr}
\pnum
\constraints
\tcode{is_same_v<remove_cvref_t<U>, in_place_t>} is \tcode{false}; and \tcode{is_same_v<remove_cvref_t<U>, unexpect_t>} is \tcode{false}; and \tcode{is_same_v<remove_cvref_t<U>, expected>} is \tcode{false}; and \tcode{is_constructible_v<T, U>} is \tcode{true}; and \tcode{remove_cvref_t<U>} is not a specialization of \tcode{unexpected}; and if \tcode{T} is \tcode{bool}, \tcode{remove_cvref_t<U>} is not a specialization of \tcode{expected}.

\pnum
\effects
Direct-non-list-initializes \tcode{val} with \tcode{std::forward<U>(v)}.

\pnum
\ensures
\tcode{has_value()} is \tcode{true}.

\pnum
\throws
Any exception thrown by the initialization of \tcode{val}.
\end{itemdescr}

\indexlibraryctor{expected}%
\begin{itemdecl}
template<class G>
  requires(!is_reference_v<E> && is_constructible_v<E, const G&>)
constexpr explicit(!is_convertible_v<const G&, E>) expected(const unexpected<G>& e);
template<class G>
  requires(!is_reference_v<E> && is_constructible_v<E, G>)
constexpr explicit(!is_convertible_v<G, E>) expected(unexpected<G>&& e);
\end{itemdecl}

\begin{itemdescr}
\pnum
\constraints
\tcode{is_constructible_v<E, const G&>} is \tcode{true}.

\pnum
\effects
Direct-non-list-initializes \tcode{unex} with \tcode{std::forward<const G&>(e.error())}.

\pnum
\ensures
\tcode{has_value()} is \tcode{false}.

\pnum
\throws
Any exception thrown by the initialization of \tcode{unex}.
\end{itemdescr}

\indexlibraryctor{expected}%
\begin{itemdecl}
template<class G>
  requires(is_reference_v<E> && is_reference_v<G> && is_constructible_v<E, G> &&
           !reference_constructs_from_temporary_v<E, G>)
constexpr explicit(!is_convertible_v<G, E>) expected(const unexpected<G>& e) noexcept;
template<class G>
  requires(is_reference_v<E> && is_reference_v<G> && is_constructible_v<E, G> &&
           !reference_constructs_from_temporary_v<E, G>)
constexpr explicit(!is_convertible_v<G, E>) expected(unexpected<G>&& e) noexcept;
\end{itemdecl}

\begin{itemdescr}
\pnum
\constraints
\tcode{is_reference_v<G>} is \tcode{true}; and \tcode{is_constructible_v<E, G>} is \tcode{true}; and \tcode{reference_constructs_from_temporary_v<E, G>} is \tcode{false}.

\pnum
\effects
Initializes \tcode{unex} with \tcode{e.error()}.

\pnum
\ensures
\tcode{has_value()} is \tcode{false}.

\pnum
\remarks
This constructor never throws: the referent is bound, not copied.
\end{itemdescr}

\indexlibraryctor{expected}%
\begin{itemdecl}
template<class G>
  requires(is_reference_v<E> && !is_reference_v<G>)
constexpr expected(const unexpected<G>&);
\end{itemdecl}

\begin{itemdescr}
\pnum
\remarks
When \tcode{E} is a reference type, an overload taking \tcode{unexpected<G>} for a non-reference \tcode{G} is defined as deleted: the referent would live inside the (possibly temporary) source \tcode{unexpected<G>} object, and binding \tcode{E&} to it would dangle. Use an \tcode{unexpected<E&>} holding an external object instead.
\end{itemdescr}

\indexlibraryctor{expected}%
\begin{itemdecl}
template<class... Args>
  requires is_constructible_v<T, Args...>
constexpr explicit expected(in_place_t, Args&&... args);
\end{itemdecl}

\begin{itemdescr}
\pnum
\constraints
\tcode{is_constructible_v<T, Args...>} is \tcode{true}.

\pnum
\effects
Direct-non-list-initializes \tcode{val} with \tcode{std::forward<Args>(args)...}.

\pnum
\ensures
\tcode{has_value()} is \tcode{true}.

\pnum
\throws
Any exception thrown by the initialization of \tcode{val}.
\end{itemdescr}

\indexlibraryctor{expected}%
\begin{itemdecl}
template<class U, class... Args>
  requires is_constructible_v<T, initializer_list<U>&, Args...>
constexpr explicit expected(in_place_t, initializer_list<U> il, Args&&... args);
\end{itemdecl}

\begin{itemdescr}
\pnum
\constraints
\tcode{is_constructible_v<T, initializer_list<U>&, Args...>} is \tcode{true}.

\pnum
\effects
Direct-non-list-initializes \tcode{val} with \tcode{il, std::forward<Args>(args)...}.

\pnum
\ensures
\tcode{has_value()} is \tcode{true}.

\pnum
\throws
Any exception thrown by the initialization of \tcode{val}.
\end{itemdescr}

\indexlibraryctor{expected}%
\begin{itemdecl}
template<class... Args>
  requires(is_constructible_v<E, Args...> && !@\exposidnc{unexpect-dangles-v}@<E, Args...>)
constexpr explicit expected(unexpect_t, Args&&... args);
\end{itemdecl}

\begin{itemdescr}
\pnum
\constraints
\tcode{is_constructible_v<E, Args...>} is \tcode{true}.

\pnum
\effects
Direct-non-list-initializes \tcode{unex} with \tcode{std::forward<Args>(args)...}.

\pnum
\ensures
\tcode{has_value()} is \tcode{false}.

\pnum
\throws
Any exception thrown by the initialization of \tcode{unex}.
\end{itemdescr}

\indexlibraryctor{expected}%
\begin{itemdecl}
template<class... Args>
  requires(@\exposidnc{unexpect-dangles-v}@<E, Args...>)
constexpr expected(unexpect_t, Args&&...);
\end{itemdecl}

\begin{itemdescr}
\pnum
\remarks
When \tcode{E} is a reference type, an overload with the same parameter types is defined as deleted if the single argument would bind \tcode{E&} to a temporary, or if it is otherwise not usable to construct \tcode{E}.
\end{itemdescr}

\indexlibraryctor{expected}%
\begin{itemdecl}
template<class U, class... Args>
  requires(!is_reference_v<E> && is_constructible_v<E, initializer_list<U>&, Args...>)
constexpr explicit expected(unexpect_t, initializer_list<U> il, Args&&... args);
\end{itemdecl}

\begin{itemdescr}
\pnum
\constraints
\tcode{is_reference_v<E>} is \tcode{false}, and \tcode{is_constructible_v<E, initializer_list<U>&, Args...>} is \tcode{true}.

\pnum
\effects
Direct-non-list-initializes \tcode{unex} with \tcode{il, std::forward<Args>(args)...}.

\pnum
\ensures
\tcode{has_value()} is \tcode{false}.

\pnum
\throws
Any exception thrown by the initialization of \tcode{unex}.

\pnum
\remarks
An overload with the same parameter types is defined as deleted when \tcode{E} is an lvalue reference type. An initializer list cannot provide the required long-lived error referent.
\end{itemdescr}

\indexlibraryctor{expected}%
\begin{itemdecl}
template<class U, class... Args>
  requires is_reference_v<E>
constexpr expected(unexpect_t, initializer_list<U>, Args&&...);
\end{itemdecl}

\begin{itemdescr}
\pnum
\remarks
An overload with the same parameter types is defined as deleted when \tcode{E} is an lvalue reference type. An initializer list cannot provide the required long-lived error referent.
\end{itemdescr}

\rSec3[expected.object.dtor]{Destructor}

\indexlibrarydtor{expected}%
\begin{itemdecl}
constexpr ~expected();
\end{itemdecl}

\begin{itemdescr}
\pnum
\constraints
\tcode{is_trivially_destructible_v<T>} is \tcode{true} and \tcode{is_trivially_destructible_v<E>} is \tcode{true}.

\pnum
\effects
None: \tcode{val} or \tcode{unex} (whichever is active) has a trivial destructor.

\pnum
\remarks
This destructor is trivial.
\end{itemdescr}

\indexlibrarydtor{expected}%
\begin{itemdecl}
constexpr ~expected();
\end{itemdecl}

\begin{itemdescr}
\pnum
\constraints
\tcode{(is_trivially_destructible_v<T> && is_trivially_destructible_v<E>)} is \tcode{false}.

\pnum
\effects
If \tcode{has_value()} is \tcode{true}, destroys \tcode{val}, otherwise destroys \tcode{unex}.
\end{itemdescr}

\rSec3[expected.object.assign]{Assignment}

\indexlibrarymember{operator=}{expected}%
\begin{itemdecl}
constexpr expected& operator=(const expected&);
\end{itemdecl}

\begin{itemdescr}
\pnum
\constraints
\begin{itemize}
\item \tcode{is_trivially_copy_constructible_v<T>} is \tcode{true},
\item \tcode{is_trivially_copy_assignable_v<T>} is \tcode{true},
\item \tcode{is_trivially_destructible_v<T>} is \tcode{true},
\item \tcode{is_trivially_copy_constructible_v<E>} is \tcode{true},
\item \tcode{is_trivially_copy_assignable_v<E>} is \tcode{true},
\item \tcode{is_trivially_destructible_v<E>} is \tcode{true}.
\end{itemize}

\pnum
\effects
Trivially copies \tcode{rhs}'s active member into \tcode{*this}.

\pnum
\returns
\tcode{*this}.

\pnum
\remarks
This operator is trivial.
\end{itemdescr}

\indexlibrarymember{operator=}{expected}%
\begin{itemdecl}
constexpr expected& operator=(const expected& rhs) noexcept(
    is_nothrow_copy_constructible_v<T> && is_nothrow_copy_assignable_v<T> &&
    is_nothrow_copy_constructible_v<E> && is_nothrow_copy_assignable_v<E>);
\end{itemdecl}

\begin{itemdescr}
\pnum
\constraints
\begin{itemize}
\item \tcode{is_copy_constructible_v<T>} is \tcode{true},
\item \tcode{is_copy_assignable_v<T>} is \tcode{true},
\item \tcode{is_reference_v<E> || (is_copy_constructible_v<E> && is_copy_assignable_v<E>)} is \tcode{true},
\item \tcode{is_nothrow_move_constructible_v<T> || is_nothrow_move_constructible_v<E>} is \tcode{true},
\item \tcode{(is_trivially_copy_constructible_v<T> && is_trivially_copy_assignable_v<T> &&
                           is_trivially_destructible_v<T> && is_trivially_copy_constructible_v<E> &&
                           is_trivially_copy_assignable_v<E> && is_trivially_destructible_v<E>)} is \tcode{false}.
\end{itemize}

\pnum
\effects
If \tcode{this->has_value() && rhs.has_value()}, equivalent to \tcode{val = *rhs}. Otherwise, if \tcode{this->has_value()}, equivalent to \tcode{reinit-expected(unex, val, rhs.error())}. Otherwise, if \tcode{rhs.has_value()}, equivalent to \tcode{reinit-expected(val, unex, *rhs)}. Otherwise, equivalent to \tcode{unex = rhs.unex}. Then, if no exception was thrown, equivalent to: \tcode{has_val = rhs.has_value(); return *this;} When \tcode{E} is an lvalue reference type, each of the cases above that initializes or assigns \tcode{unex} rebinds it: \tcode{unex} comes to refer to the same object as \tcode{rhs}'s error. No previously or subsequently referenced object is assigned through.

\pnum
\returns
\tcode{*this}.

\pnum
\remarks
This operator is defined as deleted unless \tcode{is_copy_assignable_v<T>} is \tcode{true} and \tcode{is_copy_constructible_v<T>} is \tcode{true} and \tcode{is_copy_assignable_v<E>} is \tcode{true} or \tcode{is_reference_v<E>} is \tcode{true} and \tcode{is_copy_constructible_v<E>} is \tcode{true} or \tcode{is_reference_v<E>} is \tcode{true} and \tcode{is_nothrow_move_constructible_v<T> || is_nothrow_move_constructible_v<E>} is \tcode{true}. This operator is trivial if \tcode{is_trivially_copy_constructible_v<T>}, \tcode{is_trivially_copy_assignable_v<T>}, \tcode{is_trivially_destructible_v<T>}, \tcode{is_trivially_copy_constructible_v<E>}, \tcode{is_trivially_copy_assignable_v<E>}, and \tcode{is_trivially_destructible_v<E>} are all \tcode{true}.
\end{itemdescr}

\indexlibrarymember{operator=}{expected}%
\begin{itemdecl}
constexpr expected& operator=(expected&&) noexcept;
\end{itemdecl}

\begin{itemdescr}
\pnum
\constraints
\begin{itemize}
\item \tcode{is_trivially_move_constructible_v<T>} is \tcode{true},
\item \tcode{is_trivially_move_assignable_v<T>} is \tcode{true},
\item \tcode{is_trivially_destructible_v<T>} is \tcode{true},
\item \tcode{is_trivially_move_constructible_v<E>} is \tcode{true},
\item \tcode{is_trivially_move_assignable_v<E>} is \tcode{true},
\item \tcode{is_trivially_destructible_v<E>} is \tcode{true}.
\end{itemize}

\pnum
\effects
Trivially moves \tcode{rhs}'s active member into \tcode{*this}.

\pnum
\returns
\tcode{*this}.

\pnum
\remarks
This operator is trivial.
\end{itemdescr}

\indexlibrarymember{operator=}{expected}%
\begin{itemdecl}
constexpr expected& operator=(expected&& rhs) noexcept(
    is_nothrow_move_constructible_v<T> && is_nothrow_move_assignable_v<T> &&
    is_nothrow_move_constructible_v<E> && is_nothrow_move_assignable_v<E>);
\end{itemdecl}

\begin{itemdescr}
\pnum
\constraints
\begin{itemize}
\item \tcode{is_move_constructible_v<T>} is \tcode{true},
\item \tcode{is_move_assignable_v<T>} is \tcode{true},
\item \tcode{is_reference_v<E> || (is_move_constructible_v<E> && is_move_assignable_v<E>)} is \tcode{true},
\item \tcode{is_nothrow_move_constructible_v<T> || is_nothrow_move_constructible_v<E>} is \tcode{true},
\item \tcode{(is_trivially_move_constructible_v<T> && is_trivially_move_assignable_v<T> &&
                           is_trivially_destructible_v<T> && is_trivially_move_constructible_v<E> &&
                           is_trivially_move_assignable_v<E> && is_trivially_destructible_v<E>)} is \tcode{false}.
\end{itemize}

\pnum
\effects
If \tcode{this->has_value() && rhs.has_value()}, equivalent to \tcode{val = std::move(*rhs)}. Otherwise, if \tcode{this->has_value()}, equivalent to \tcode{reinit-expected(unex, val, std::move(rhs.error()))}. Otherwise, if \tcode{rhs.has_value()}, equivalent to \tcode{reinit-expected(val, unex, std::move(*rhs))}. Otherwise, equivalent to \tcode{unex = std::move(rhs.unex)}. Then, if no exception was thrown, equivalent to: \tcode{has_val = rhs.has_value(); return *this;} When \tcode{E} is an lvalue reference type, each of the cases above that initializes or assigns \tcode{unex} rebinds it: \tcode{unex} comes to refer to the same object as \tcode{rhs}'s error. No previously or subsequently referenced object is assigned through.

\pnum
\returns
\tcode{*this}.

\pnum
\remarks
The exception specification is equivalent to \tcode{is_nothrow_move_assignable_v<T> && is_nothrow_move_constructible_v<T> && is_nothrow_move_assignable_v<E> && is_nothrow_move_constructible_v<E>}. This operator is trivial if \tcode{is_trivially_move_constructible_v<T>}, \tcode{is_trivially_move_assignable_v<T>}, \tcode{is_trivially_destructible_v<T>}, \tcode{is_trivially_move_constructible_v<E>}, \tcode{is_trivially_move_assignable_v<E>}, and \tcode{is_trivially_destructible_v<E>} are all \tcode{true}.
\end{itemdescr}

\indexlibrarymember{operator=}{expected}%
\begin{itemdecl}
template<class U = remove_cv_t<T>>
  requires(!is_same_v<expected, remove_cvref_t<U>> &&
           !@\exposidnc{is-unexpected-specialization}@<remove_cvref_t<U>>::value &&
           is_constructible_v<T, U> && is_assignable_v<T&, U> &&
           (is_nothrow_constructible_v<T, U> || is_nothrow_move_constructible_v<T> ||
            is_nothrow_move_constructible_v<E>))
constexpr expected& operator=(U&& v);
\end{itemdecl}

\begin{itemdescr}
\pnum
\constraints
\tcode{is_same_v<expected, remove_cvref_t<U>>} is \tcode{false}; and \tcode{remove_cvref_t<U>} is not a specialization of \tcode{unexpected}; and \tcode{is_constructible_v<T, U>} is \tcode{true}; and \tcode{is_assignable_v<T&, U>} is \tcode{true}; and \tcode{is_nothrow_constructible_v<T, U> || is_nothrow_move_constructible_v<T> || is_nothrow_move_constructible_v<E>} is \tcode{true}.

\pnum
\effects
If \tcode{has_value()} is \tcode{true}, equivalent to: \tcode{val = std::forward<U>(v);} Otherwise, equivalent to: \tcode{reinit-expected(val, unex, std::forward<U>(v)); has_val = true;}

\pnum
\returns
\tcode{*this}.
\end{itemdescr}

\indexlibrarymember{operator=}{expected}%
\begin{itemdecl}
template<class G>
  requires(!is_reference_v<E> && is_constructible_v<E, const G&> &&
           is_assignable_v<E&, const G&> &&
           (is_nothrow_constructible_v<E, const G&> ||
            is_nothrow_move_constructible_v<T> || is_nothrow_move_constructible_v<E>))
constexpr expected& operator=(const unexpected<G>& e);
template<class G>
  requires(!is_reference_v<E> && is_constructible_v<E, G> && is_assignable_v<E&, G> &&
           (is_nothrow_constructible_v<E, G> || is_nothrow_move_constructible_v<T> ||
            is_nothrow_move_constructible_v<E>))
constexpr expected& operator=(unexpected<G>&& e);
\end{itemdecl}

\begin{itemdescr}
\pnum
\constraints
\tcode{is_constructible_v<E, const G&>} is \tcode{true}; and \tcode{is_assignable_v<E&, const G&>} is \tcode{true}; and \tcode{is_nothrow_constructible_v<E, const G&> || is_nothrow_move_constructible_v<T> || is_nothrow_move_constructible_v<E>} is \tcode{true}.

\pnum
\effects
If \tcode{has_value()} is \tcode{true}, equivalent to: \tcode{reinit-expected(unex, val, e.error()); has_val = false;} Otherwise, equivalent to: \tcode{unex = unexpected<E>(e.error());}

\pnum
\returns
\tcode{*this}.
\end{itemdescr}

\indexlibrarymember{operator=}{expected}%
\begin{itemdecl}
template<class G>
  requires(is_reference_v<E> && is_reference_v<G> && is_constructible_v<E, G> &&
           !reference_constructs_from_temporary_v<E, G>)
constexpr expected& operator=(const unexpected<G>& e);
template<class G>
  requires(is_reference_v<E> && is_reference_v<G> && is_constructible_v<E, G> &&
           !reference_constructs_from_temporary_v<E, G>)
constexpr expected& operator=(unexpected<G>&& e);
\end{itemdecl}

\begin{itemdescr}
\pnum
\constraints
\tcode{is_reference_v<G>} is \tcode{true}; and \tcode{is_constructible_v<E, G>} is \tcode{true}; and \tcode{reference_constructs_from_temporary_v<E, G>} is \tcode{false}.

\pnum
\effects
Rebinds \tcode{unex} to refer to the same object as \tcode{e.error()}, destroying \tcode{val} first if \tcode{has_value()} is \tcode{true}.

\pnum
\ensures
\tcode{has_value()} is \tcode{false}.

\pnum
\returns
\tcode{*this}.

\pnum
\remarks
This operator never throws: the referent is bound, not copied.
\end{itemdescr}

\indexlibrarymember{operator=}{expected}%
\begin{itemdecl}
template<class G>
  requires(is_reference_v<E> && !is_reference_v<G>)
constexpr expected& operator=(const unexpected<G>&);
\end{itemdecl}

\begin{itemdescr}
\pnum
\remarks
When \tcode{E} is a reference type, an overload taking \tcode{unexpected<G>} for a non-reference \tcode{G} is defined as deleted: it would rebind \tcode{unex} to \tcode{unexpected<G>}'s temporary storage.
\end{itemdescr}

\indexlibrarymember{emplace}{expected}%
\begin{itemdecl}
template<class... Args>
  requires is_nothrow_constructible_v<T, Args...>
constexpr T& emplace(Args&&... args) noexcept;
\end{itemdecl}

\begin{itemdescr}
\pnum
\constraints
\tcode{is_nothrow_constructible_v<T, Args...>} is \tcode{true}.

\pnum
\effects
Equivalent to: \tcode{if (has_value()) { destroy_at(addressof(val)); } else { destroy_at(addressof(unex)); has_val = true; } return *construct_at(addressof(val), std::forward<Args>(args)...);}
\end{itemdescr}

\indexlibrarymember{emplace}{expected}%
\begin{itemdecl}
template<class U, class... Args>
  requires is_nothrow_constructible_v<T, initializer_list<U>&, Args...>
constexpr T& emplace(initializer_list<U> il, Args&&... args) noexcept;
\end{itemdecl}

\begin{itemdescr}
\pnum
\constraints
\tcode{is_nothrow_constructible_v<T, initializer_list<U>&, Args...>} is \tcode{true}.

\pnum
\effects
Equivalent to: \tcode{if (has_value()) { destroy_at(addressof(val)); } else { destroy_at(addressof(unex)); has_val = true; } return *construct_at(addressof(val), il, std::forward<Args>(args)...);}
\end{itemdescr}

\rSec3[expected.object.swap]{Swap}

\indexlibrarymember{swap}{expected}%
\begin{itemdecl}
constexpr void swap(expected& rhs) noexcept(is_nothrow_move_constructible_v<T> &&
                                            is_nothrow_swappable_v<T> &&
                                            is_nothrow_move_constructible_v<E> &&
                                            (is_reference_v<E> ||
                                             is_nothrow_swappable_v<E>));
\end{itemdecl}

\begin{itemdescr}
\pnum
\constraints
\begin{itemize}
\item \tcode{is_swappable_v<T>} is \tcode{true},
\item \tcode{is_reference_v<E> || is_swappable_v<E>} is \tcode{true},
\item \tcode{is_move_constructible_v<T>} is \tcode{true},
\item \tcode{is_move_constructible_v<E>} is \tcode{true},
\item \tcode{is_nothrow_move_constructible_v<T> || is_nothrow_move_constructible_v<E>} is \tcode{true}.
\end{itemize}

\pnum
\effects
If \tcode{this->has_value()} and \tcode{rhs.has_value()}, equivalent to \tcode{using std::swap; swap(val, rhs.val);}. If neither \tcode{*this} nor \tcode{rhs} contains a value, equivalent to \tcode{using std::swap; swap(unex, rhs.unex);}. If \tcode{rhs.has_value()} is \tcode{false} and \tcode{this->has_value()} is \tcode{true}, exchanges the value and error between \tcode{*this} and \tcode{rhs} (moving through a temporary so a failed move leaves both objects unchanged), leaving \tcode{has_value()} \tcode{false} and \tcode{rhs.has_value()} \tcode{true}. If \tcode{rhs.has_value()} is \tcode{true} and \tcode{this->has_value()} is \tcode{false}, equivalent to \tcode{rhs.swap(*this)}.

\pnum
\throws
Any exception thrown by the expressions in the Effects.

\pnum
\remarks
The exception specification is equivalent to \tcode{is_nothrow_move_constructible_v<T> && is_nothrow_swappable_v<T> && is_nothrow_move_constructible_v<E> && (is_reference_v<E> || is_nothrow_swappable_v<E>)}.
\end{itemdescr}

\indexlibraryglobal{swap}%
\begin{itemdecl}
friend constexpr void swap(expected& x, expected& y) noexcept(noexcept(x.swap(y)));
\end{itemdecl}

\begin{itemdescr}
\pnum
\constraints
\begin{itemize}
\item \tcode{is_swappable_v<T>} is \tcode{true},
\item \tcode{is_reference_v<E> || is_swappable_v<E>} is \tcode{true},
\item \tcode{is_move_constructible_v<T>} is \tcode{true},
\item \tcode{is_move_constructible_v<E>} is \tcode{true},
\item \tcode{is_nothrow_move_constructible_v<T> || is_nothrow_move_constructible_v<E>} is \tcode{true}.
\end{itemize}

\pnum
\effects
Equivalent to \tcode{x.swap(y)}.
\end{itemdescr}

\rSec3[expected.object.obs]{Observers}

\indexlibrarymember{operator->}{expected}%
\begin{itemdecl}
constexpr const T* operator->() const noexcept;
constexpr T* operator->() noexcept;
\end{itemdecl}

\begin{itemdescr}
\pnum
\hardexpects
\tcode{has_value()} is \tcode{true}.

\pnum
\returns
\tcode{addressof(val)}.
\end{itemdescr}

\indexlibrarymember{operator*}{expected}%
\begin{itemdecl}
constexpr const T& operator*() const& noexcept;
constexpr T& operator*() & noexcept;
\end{itemdecl}

\begin{itemdescr}
\pnum
\hardexpects
\tcode{has_value()} is \tcode{true}.

\pnum
\returns
\tcode{val}.
\end{itemdescr}

\indexlibrarymember{operator*}{expected}%
\begin{itemdecl}
constexpr const T&& operator*() const&& noexcept;
constexpr T&& operator*() && noexcept;
\end{itemdecl}

\begin{itemdescr}
\pnum
\hardexpects
\tcode{has_value()} is \tcode{true}.

\pnum
\returns
\tcode{std::move(val)}.
\end{itemdescr}

\indexlibrarymember{operator bool}{expected}%
\indexlibrarymember{has_value}{expected}%
\begin{itemdecl}
constexpr explicit operator bool() const noexcept;
constexpr bool has_value() const noexcept;
\end{itemdecl}

\begin{itemdescr}
\pnum
\returns
\tcode{has_val}.
\end{itemdescr}

\indexlibrarymember{value}{expected}%
\begin{itemdecl}
constexpr const T& value() const&;
constexpr T& value() &;
\end{itemdecl}

\begin{itemdescr}
\pnum
\mandates
\tcode{is_copy_constructible_v<E>} is \tcode{true}.

\pnum
\returns
\tcode{val}, if \tcode{has_value()} is \tcode{true}.

\pnum
\throws
\tcode{bad_expected_access(as_const(error()))} if \tcode{has_value()} is \tcode{false}.
\end{itemdescr}

\indexlibrarymember{value}{expected}%
\begin{itemdecl}
constexpr const T&& value() const&&;
constexpr T&& value() &&;
\end{itemdecl}

\begin{itemdescr}
\pnum
\mandates
\tcode{is_copy_constructible_v<E>} is \tcode{true} and \tcode{is_constructible_v<E, decltype(std::move(error()))>} is \tcode{true}.

\pnum
\returns
\tcode{std::move(val)}, if \tcode{has_value()} is \tcode{true}.

\pnum
\throws
\tcode{bad_expected_access(std::move(error()))} if \tcode{has_value()} is \tcode{false}.
\end{itemdescr}

\indexlibrarymember{error}{expected}%
\begin{itemdecl}
constexpr const E& error() const& noexcept;
constexpr E& error() & noexcept;
\end{itemdecl}

\begin{itemdescr}
\pnum
\hardexpects
\tcode{has_value()} is \tcode{false}.

\pnum
\returns
\tcode{unex.error()}.
\end{itemdescr}

\indexlibrarymember{error}{expected}%
\begin{itemdecl}
constexpr const E&& error() const&& noexcept;
constexpr E&& error() && noexcept;
\end{itemdecl}

\begin{itemdescr}
\pnum
\hardexpects
\tcode{has_value()} is \tcode{false}.

\pnum
\returns
\tcode{std::move(unex).error()}.
\end{itemdescr}

\indexlibrarymember{value_or}{expected}%
\begin{itemdecl}
template<class U = remove_cv_t<T>> constexpr T value_or(U&& def) const&;
\end{itemdecl}

\begin{itemdescr}
\pnum
\mandates
\tcode{is_copy_constructible_v<T>} is \tcode{true} and \tcode{is_convertible_v<U, T>} is \tcode{true}.

\pnum
\returns
\tcode{has_value() ? **this : static_cast<T>(std::forward<U>(def))}.
\end{itemdescr}

\indexlibrarymember{value_or}{expected}%
\begin{itemdecl}
template<class U = remove_cv_t<T>> constexpr T value_or(U&& def) &&;
\end{itemdecl}

\begin{itemdescr}
\pnum
\mandates
\tcode{is_move_constructible_v<T>} is \tcode{true} and \tcode{is_convertible_v<U, T>} is \tcode{true}.

\pnum
\returns
\tcode{has_value() ? std::move(**this) : static_cast<T>(std::forward<U>(def))}.
\end{itemdescr}

\indexlibrarymember{error_or}{expected}%
\begin{itemdecl}
template<class G = @\exposidnc{error-value-type}@>
  requires(is_copy_constructible_v<@\exposidnc{error-value-type}@> &&
           is_convertible_v<G, @\exposidnc{error-value-type}@>)
constexpr @\exposidnc{error-value-type}@ error_or(G&& def) const&;
\end{itemdecl}

\begin{itemdescr}
\pnum
\mandates
\tcode{is_copy_constructible_v<error_value_type>} is \tcode{true} and \tcode{is_convertible_v<G, error_value_type>} is \tcode{true}.

\pnum
\returns
\tcode{std::forward<G>(def)} if \tcode{has_value()} is \tcode{true}, \tcode{error()} otherwise.
\end{itemdescr}

\indexlibrarymember{error_or}{expected}%
\begin{itemdecl}
template<class G = @\exposidnc{error-value-type}@>
  requires(is_move_constructible_v<@\exposidnc{error-value-type}@> &&
           is_convertible_v<G, @\exposidnc{error-value-type}@>)
constexpr @\exposidnc{error-value-type}@ error_or(G&& def) &&;
\end{itemdecl}

\begin{itemdescr}
\pnum
\mandates
\tcode{is_move_constructible_v<error_value_type>} is \tcode{true} and \tcode{is_convertible_v<G, error_value_type>} is \tcode{true}.

\pnum
\returns
\tcode{std::forward<G>(def)} if \tcode{has_value()} is \tcode{true}, \tcode{std::move(error())} otherwise.
\end{itemdescr}

\rSec3[expected.object.monadic]{Monadic operations}

\indexlibrarymember{and_then}{expected}%
\begin{itemdecl}
template<class F>
  requires is_constructible_v<E, E&>
constexpr auto and_then(F&& f) &;
template<class F>
  requires is_constructible_v<E, const E&>
constexpr auto and_then(F&& f) const&;
\end{itemdecl}

\begin{itemdescr}
\pnum
\constraints
\tcode{is_constructible_v<E, decltype(error())>} is \tcode{true}.

\pnum
\mandates
\tcode{remove_cvref_t<invoke_result_t<F, decltype((val))>>} is a specialization of \tcode{expected} and its \tcode{error_type} is the same type as \tcode{E}.

\pnum
\effects
Equivalent to: \tcode{if (has_value()) return invoke(std::forward<F>(f), val); else return U(unexpect, error());} where \tcode{U} is \tcode{remove_cvref_t<invoke_result_t<F, decltype((val))>>}.
\end{itemdescr}

\indexlibrarymember{and_then}{expected}%
\begin{itemdecl}
template<class F>
  requires is_constructible_v<E, E&&>
constexpr auto and_then(F&& f) &&;
template<class F>
  requires is_constructible_v<E, const E&&>
constexpr auto and_then(F&& f) const&&;
\end{itemdecl}

\begin{itemdescr}
\pnum
\constraints
\tcode{is_constructible_v<E, decltype(std::move(error()))>} is \tcode{true}.

\pnum
\mandates
\tcode{remove_cvref_t<invoke_result_t<F, decltype(std::move(val))>>} is a specialization of \tcode{expected} and its \tcode{error_type} is the same type as \tcode{E}.

\pnum
\effects
Equivalent to: \tcode{if (has_value()) return invoke(std::forward<F>(f), std::move(val)); else return U(unexpect, std::move(error()));} where \tcode{U} is \tcode{remove_cvref_t<invoke_result_t<F, decltype(std::move(val))>>}.
\end{itemdescr}

\indexlibrarymember{or_else}{expected}%
\begin{itemdecl}
template<class F>
  requires is_constructible_v<T, T&>
constexpr auto or_else(F&& f) &;
template<class F>
  requires is_constructible_v<T, const T&>
constexpr auto or_else(F&& f) const&;
\end{itemdecl}

\begin{itemdescr}
\pnum
\constraints
\tcode{is_constructible_v<T, decltype((val))>} is \tcode{true}.

\pnum
\mandates
\tcode{remove_cvref_t<invoke_result_t<F, decltype(error())>>} is a specialization of \tcode{expected} and its \tcode{value_type} is the same type as \tcode{T}.

\pnum
\effects
Equivalent to: \tcode{if (has_value()) return G(in_place, val); else return invoke(std::forward<F>(f), error());} where \tcode{G} is \tcode{remove_cvref_t<invoke_result_t<F, decltype(error())>>}.
\end{itemdescr}

\indexlibrarymember{or_else}{expected}%
\begin{itemdecl}
template<class F>
  requires is_constructible_v<T, T&&>
constexpr auto or_else(F&& f) &&;
template<class F>
  requires is_constructible_v<T, const T&&>
constexpr auto or_else(F&& f) const&&;
\end{itemdecl}

\begin{itemdescr}
\pnum
\constraints
\tcode{is_constructible_v<T, decltype(std::move(val))>} is \tcode{true}.

\pnum
\mandates
\tcode{remove_cvref_t<invoke_result_t<F, decltype(std::move(error()))>>} is a specialization of \tcode{expected} and its \tcode{value_type} is the same type as \tcode{T}.

\pnum
\effects
Equivalent to: \tcode{if (has_value()) return G(in_place, std::move(val)); else return invoke(std::forward<F>(f), std::move(error()));} where \tcode{G} is \tcode{remove_cvref_t<invoke_result_t<F, decltype(std::move(error()))>>}.
\end{itemdescr}

\indexlibrarymember{transform}{expected}%
\begin{itemdecl}
template<class F>
  requires is_constructible_v<E, E&>
constexpr auto transform(F&& f) &;
template<class F>
  requires is_constructible_v<E, const E&>
constexpr auto transform(F&& f) const&;
\end{itemdecl}

\begin{itemdescr}
\pnum
\constraints
\tcode{is_constructible_v<E, decltype(error())>} is \tcode{true}.

\pnum
\effects
Equivalent to: \tcode{if (!has_value()) return U(unexpect, error()); else return expected<U2, E>(in_place, invoke(std::forward<F>(f), val));} where \tcode{U2} is \tcode{remove_cv_t<invoke_result_t<F, decltype((val))>>} and \tcode{U} is \tcode{expected<U2, E>}.
\end{itemdescr}

\indexlibrarymember{transform}{expected}%
\begin{itemdecl}
template<class F>
  requires is_constructible_v<E, E&&>
constexpr auto transform(F&& f) &&;
template<class F>
  requires is_constructible_v<E, const E&&>
constexpr auto transform(F&& f) const&&;
\end{itemdecl}

\begin{itemdescr}
\pnum
\constraints
\tcode{is_constructible_v<E, decltype(std::move(error()))>} is \tcode{true}.

\pnum
\effects
Equivalent to: \tcode{if (!has_value()) return U(unexpect, std::move(error())); else return expected<U2, E>(in_place, invoke(std::forward<F>(f), std::move(val)));} where \tcode{U2} is \tcode{remove_cv_t<invoke_result_t<F, decltype(std::move(val))>>} and \tcode{U} is \tcode{expected<U2, E>}.
\end{itemdescr}

\indexlibrarymember{transform_error}{expected}%
\begin{itemdecl}
template<class F>
  requires is_constructible_v<T, T&>
constexpr auto transform_error(F&& f) &;
template<class F>
  requires is_constructible_v<T, const T&>
constexpr auto transform_error(F&& f) const&;
\end{itemdecl}

\begin{itemdescr}
\pnum
\constraints
\tcode{is_constructible_v<T, decltype((val))>} is \tcode{true}.

\pnum
\mandates
\begin{itemize}
\item \tcode{is_object_v<G>} is \tcode{true},
\item \tcode{is_array_v<G>} is \tcode{false},
\item \tcode{is_same_v<G, remove_cv_t<G>>} is \tcode{true},
\item \tcode{\exposid{is-unexpected-specialization}<G>::value} is \tcode{false}.
\end{itemize}

\pnum
\effects
Equivalent to: \tcode{if (has_value()) return G(in_place, val); else return expected<T, G2>(unexpect, invoke(std::forward<F>(f), error()));} where \tcode{G2} is \tcode{remove_cv_t<invoke_result_t<F, decltype(error())>>} and \tcode{G} is \tcode{expected<T, G2>}.
\end{itemdescr}

\indexlibrarymember{transform_error}{expected}%
\begin{itemdecl}
template<class F>
  requires is_constructible_v<T, T&&>
constexpr auto transform_error(F&& f) &&;
template<class F>
  requires is_constructible_v<T, const T&&>
constexpr auto transform_error(F&& f) const&&;
\end{itemdecl}

\begin{itemdescr}
\pnum
\constraints
\tcode{is_constructible_v<T, decltype(std::move(val))>} is \tcode{true}.

\pnum
\mandates
\begin{itemize}
\item \tcode{is_object_v<G>} is \tcode{true},
\item \tcode{is_array_v<G>} is \tcode{false},
\item \tcode{is_same_v<G, remove_cv_t<G>>} is \tcode{true},
\item \tcode{\exposid{is-unexpected-specialization}<G>::value} is \tcode{false}.
\end{itemize}

\pnum
\effects
Equivalent to: \tcode{if (has_value()) return G(in_place, std::move(val)); else return expected<T, G2>(unexpect, invoke(std::forward<F>(f), std::move(error())));} where \tcode{G2} is \tcode{remove_cv_t<invoke_result_t<F, decltype(std::move(error()))>>} and \tcode{G} is \tcode{expected<T, G2>}.
\end{itemdescr}

\rSec3[expected.object.eq]{Equality operators}

\indexlibraryglobal{operator==}%
\begin{itemdecl}
template<class T2, class E2>
  requires(!is_void_v<T2>)
friend constexpr bool operator==(const expected& x, const expected<T2, E2>& y);
\end{itemdecl}

\begin{itemdescr}
\pnum
\mandates
\tcode{!is_void_v<T2>} is \tcode{true}. The expression \tcode{*x == *y} is well-formed and its result is convertible to \tcode{bool}. The expression \tcode{x.error() == y.error()} is well-formed and its result is convertible to \tcode{bool}.

\pnum
\returns
If \tcode{x.has_value() != y.has_value()}, \tcode{false}; otherwise, if \tcode{x.has_value()} is \tcode{true}, \tcode{*x == *y}; otherwise \tcode{x.error() == y.error()}.
\end{itemdescr}

\indexlibraryglobal{operator==}%
\begin{itemdecl}
template<class T2>
  requires(!@\exposidnc{is-expected-specialization}@<T2>::value)
friend constexpr bool operator==(const expected& x, const T2& val);
\end{itemdecl}

\begin{itemdescr}
\pnum
\mandates
\tcode{T2} is not a specialization of \tcode{expected}. The expression \tcode{*x == val} is well-formed and its result is convertible to \tcode{bool}.

\pnum
\returns
\tcode{x.has_value() && static_cast<bool>(*x == val)}.
\end{itemdescr}

\indexlibraryglobal{operator==}%
\begin{itemdecl}
template<class E2>
friend constexpr bool operator==(const expected& x, const unexpected<E2>& e);
\end{itemdecl}

\begin{itemdescr}
\pnum
\mandates
The expression \tcode{x.error() == e.error()} is well-formed and its result is convertible to \tcode{bool}.

\pnum
\returns
\tcode{!x.has_value() && static_cast<bool>(x.error() == e.error())}.
\end{itemdescr}

\rSec2[expected.void]{Partial specialization of expected for void types}

\rSec3[expected.void.general]{General}

\begin{codeblock}
template<class E>
class @\libglobal{expected}@<void, E> {
private:
  using @\exposidnc{error-value-type}@ = remove_cv_t<remove_reference_t<E>>; // exposition only

public:
  using @\libmember{value_type}{expected}@ = void;
  using @\libmember{error_type}{expected}@ = E;
  using @\libmember{unexpected_type}{expected}@ = unexpected<E>;

  template<class U> using rebind = expected<U, error_type>;

  // -------------------------------------------------------------------------
  // [expected.void.cons] Constructors
  // -------------------------------------------------------------------------

  constexpr expected() noexcept;

  // Unconstrained trivial-path candidate: see the primary template's copy
  // constructor for why (the sole declaration when E is not copy
  // constructible at all; subsumed by the non-trivial path otherwise).

  constexpr expected(const expected& rhs) noexcept(is_nothrow_copy_constructible_v<E>)
    requires(is_copy_constructible_v<E> && !is_trivially_copy_constructible_v<E>);

  // Unconstrained; no explicit noexcept — see the primary template's move
  // constructor for why.

  constexpr expected(expected&& rhs) noexcept(is_nothrow_move_constructible_v<E>)
    requires(is_move_constructible_v<E> && !is_trivially_move_constructible_v<E>);

  // Converting constructor from expected<U, G> where is_void_v<U>. Excludes U,G exactly
  // matching this class's own void,E (the real copy/move constructors already handle
  // that case) — instantiating this template for the self-referential case would
  // otherwise probe unexpected<E>'s constructibility from this very class, which some
  // standard library implementations of reference_constructs_from_temporary_v resolve
  // as a circular constraint.
  template<class U, class G>
    requires(is_void_v<U> && !is_reference_v<E> && !is_same_v<G, E> &&
             is_constructible_v<E, const G&> &&
             !is_constructible_v<unexpected<E>, expected<U, G>&> &&
             !is_constructible_v<unexpected<E>, expected<U, G>&&> &&
             !is_constructible_v<unexpected<E>, const expected<U, G>&> &&
             !is_constructible_v<unexpected<E>, const expected<U, G>&&>)
  constexpr explicit(!is_convertible_v<const G&, E>)
      expected(const expected<U, G>& rhs);

  template<class U, class G>
    requires(is_void_v<U> && !is_reference_v<E> && !is_same_v<G, E> &&
             is_constructible_v<E, G> &&
             !is_constructible_v<unexpected<E>, expected<U, G>&> &&
             !is_constructible_v<unexpected<E>, expected<U, G>&&> &&
             !is_constructible_v<unexpected<E>, const expected<U, G>&> &&
             !is_constructible_v<unexpected<E>, const expected<U, G>&&>)
  constexpr explicit(!is_convertible_v<G, E>) expected(expected<U, G>&& rhs);

  // Constructor from unexpected<G> const& / && — value-E path
  template<class G>
    requires(!is_reference_v<E> && is_constructible_v<E, const G&>)
  constexpr explicit(!is_convertible_v<const G&, E>) expected(const unexpected<G>& e);

  template<class G>
    requires(!is_reference_v<E> && is_constructible_v<E, G>)
  constexpr explicit(!is_convertible_v<G, E>) expected(unexpected<G>&& e);

  // Constructor from unexpected<G> — reference-E path. Allowed only when G is itself a
  // reference, so e.error() refers to an external object and binding E& cannot dangle.
  // No const_cast needed.
  template<class G>
    requires(is_reference_v<E> && is_reference_v<G> && is_constructible_v<E, G> &&
             !reference_constructs_from_temporary_v<E, G>)
  constexpr explicit(!is_convertible_v<G, E>) expected(const unexpected<G>& e) noexcept;

  template<class G>
    requires(is_reference_v<E> && is_reference_v<G> && is_constructible_v<E, G> &&
             !reference_constructs_from_temporary_v<E, G>)
  constexpr explicit(!is_convertible_v<G, E>) expected(unexpected<G>&& e) noexcept;

  // Deleted for reference E with value G: the referent lives inside the temporary
  // unexpected<G>, so binding E& to it would dangle once the source is destroyed. Use
  // (unexpect, lvalue) instead.
  template<class G>
    requires(is_reference_v<E> && !is_reference_v<G>)
  constexpr expected(const unexpected<G>&) = BEMAN_EXPECTED_DELETE_MSG(
      "expected<void,E&>: cannot construct from unexpected<value>; the value would "
      "dangle — use unexpected<E&>");

  template<class G>
    requires(is_reference_v<E> && !is_reference_v<G>)
  constexpr expected(unexpected<G>&&) = BEMAN_EXPECTED_DELETE_MSG(
      "expected<void,E&>: cannot construct from unexpected<value>; the value would "
      "dangle — use unexpected<E&>");

  // In-place constructor for value (no args, just marks has-value)
  constexpr explicit expected(in_place_t) noexcept;

  // In-place constructor for error
  template<class... Args>
    requires(is_constructible_v<E, Args...> && !@\exposidnc{unexpect-dangles-v}@<E, Args...>)
  constexpr explicit expected(unexpect_t, Args&&... args);

  // Deleted: single argument would bind E& to a temporary — dangling prevention
  template<class... Args>
    requires(@\exposidnc{unexpect-dangles-v}@<E, Args...>)
  constexpr expected(unexpect_t, Args&&...) =
      BEMAN_EXPECTED_DELETE_MSG("expected<void,E&>: unexpect argument would bind a "
                                "temporary that dangles; pass an lvalue reference");

  // Deleted catch-all: reference E, argument neither constructible nor a dangling case
  // (e.g. binding a non-const E& from a const lvalue).
  template<class... Args>
    requires(is_reference_v<E> && !is_constructible_v<E, Args...> &&
             !@\exposidnc{unexpect-dangles-v}@<E, Args...>)
  constexpr expected(unexpect_t, Args&&...) = BEMAN_EXPECTED_DELETE_MSG(
      "expected<void,E&>: no viable conversion from the given argument(s) to E&");

  // In-place constructor for error with initializer_list
  template<class U, class... Args>
    requires(!is_reference_v<E> && is_constructible_v<E, initializer_list<U>&, Args...>)
  constexpr explicit expected(unexpect_t, initializer_list<U> il, Args&&... args);

  template<class U, class... Args>
    requires is_reference_v<E>
  constexpr expected(unexpect_t, initializer_list<U>, Args&&...) =
      BEMAN_EXPECTED_DELETE_MSG(
          "expected<void,E&>: initializer-list error construction cannot bind a "
          "reference; pass an lvalue reference");

  // Converting constructor from expected<void, G&> — reference-E path only. G is itself
  // a reference to an external object, so binding E& to it cannot dangle regardless of
  // the source's value category, provided the reference conversion itself does not
  // materialize a temporary (e.g. a base-from-derived or qualification conversion is
  // fine; a user-defined conversion that returns by value is not). Mirrors the
  // unexpected<G> reference-E path above.
  template<class G>
    requires(is_reference_v<E> && is_convertible_v<G&, E> &&
             !reference_constructs_from_temporary_v<E, G&>)
  constexpr explicit(!is_convertible_v<G&, E>) expected(const expected<void, G&>& rhs);

  template<class G>
    requires(is_reference_v<E> && is_convertible_v<G&, E> &&
             !reference_constructs_from_temporary_v<E, G&>)
  constexpr explicit(!is_convertible_v<G&, E>) expected(expected<void, G&>&& rhs);

  // -------------------------------------------------------------------------
  // [expected.void.dtor] Destructor
  // -------------------------------------------------------------------------

  constexpr ~expected()
    requires(!is_trivially_destructible_v<E>);

  // -------------------------------------------------------------------------
  // [expected.void.assign] Assignment
  // -------------------------------------------------------------------------

  // Copy assignment (trivial path)

  // Copy assignment (non-trivial path)
  constexpr expected& operator=(const expected& rhs) noexcept(
      is_nothrow_copy_constructible_v<E> && is_nothrow_copy_assignable_v<E>)
    requires((is_reference_v<E> ||
              (is_copy_constructible_v<E> && is_copy_assignable_v<E>)) &&
             !(is_trivially_copy_constructible_v<E> &&
               is_trivially_copy_assignable_v<E> && is_trivially_destructible_v<E>));

  // Move assignment (trivial path)

  // Move assignment (non-trivial path)
  constexpr expected& operator=(expected&& rhs) noexcept(
      is_nothrow_move_constructible_v<E> && is_nothrow_move_assignable_v<E>)
    requires((is_reference_v<E> ||
              (is_move_constructible_v<E> && is_move_assignable_v<E>)) &&
             !(is_trivially_move_constructible_v<E> &&
               is_trivially_move_assignable_v<E> && is_trivially_destructible_v<E>));

  // Assignment from unexpected<G> — value-E path.
  template<class G>
    requires(!is_reference_v<E> && is_constructible_v<E, const G&> &&
             is_assignable_v<E&, const G&>)
  constexpr expected& operator=(const unexpected<G>& e);

  template<class G>
    requires(!is_reference_v<E> && is_constructible_v<E, G> && is_assignable_v<E&, G>)
  constexpr expected& operator=(unexpected<G>&& e);

  // Rebinding assignment for reference E from reference G — binds unex_ to the external
  // referent.
  template<class G>
    requires(is_reference_v<E> && is_reference_v<G> && is_constructible_v<E, G> &&
             !reference_constructs_from_temporary_v<E, G>)
  constexpr expected& operator=(const unexpected<G>& e);

  template<class G>
    requires(is_reference_v<E> && is_reference_v<G> && is_constructible_v<E, G> &&
             !reference_constructs_from_temporary_v<E, G>)
  constexpr expected& operator=(unexpected<G>&& e);

  // Deleted for reference E with value G: would bind E& to storage inside the temporary
  // unexpected.
  template<class G>
    requires(is_reference_v<E> && !is_reference_v<G>)
  constexpr expected& operator=(const unexpected<G>&) = BEMAN_EXPECTED_DELETE_MSG(
      "expected<void,E&>: cannot assign from unexpected<value>; the value would dangle "
      "— use unexpected<E&>");

  template<class G>
    requires(is_reference_v<E> && !is_reference_v<G>)
  constexpr expected& operator=(unexpected<G>&&) = BEMAN_EXPECTED_DELETE_MSG(
      "expected<void,E&>: cannot assign from unexpected<value>; the value would dangle "
      "— use unexpected<E&>");

  constexpr void emplace() noexcept;

  // -------------------------------------------------------------------------
  // [expected.void.swap] Swap
  // -------------------------------------------------------------------------

  constexpr void swap(expected& rhs) noexcept(is_nothrow_move_constructible_v<E> &&
                                              (is_reference_v<E> ||
                                               is_nothrow_swappable_v<E>))
    requires((is_reference_v<E> || is_swappable_v<E>) && is_move_constructible_v<E>);

  friend constexpr void swap(expected& x, expected& y) noexcept(noexcept(x.swap(y)))
    requires((is_reference_v<E> || is_swappable_v<E>) && is_move_constructible_v<E>);

  // -------------------------------------------------------------------------
  // [expected.void.obs] Observers
  // -------------------------------------------------------------------------

  constexpr explicit operator bool() const noexcept;
  constexpr bool has_value() const noexcept;

  constexpr void operator*() const noexcept;

  constexpr void value() const&;
  constexpr void value() &&;

  // error() — shallow const for reference E: always returns E& regardless of const on
  // expected
  constexpr const E& error() const& noexcept;
  constexpr E& error() & noexcept;
  constexpr const E&& error() const&& noexcept;
  constexpr E&& error() && noexcept;

  template<class G = @\exposidnc{error-value-type}@>
    requires(is_copy_constructible_v<remove_cv_t<remove_reference_t<E>>> &&
             is_convertible_v<G, remove_cv_t<remove_reference_t<E>>>)
  constexpr @\exposidnc{error-value-type}@ error_or(G&& def) const&;

  template<class G = @\exposidnc{error-value-type}@>
    requires(is_move_constructible_v<remove_cv_t<remove_reference_t<E>>> &&
             is_convertible_v<G, remove_cv_t<remove_reference_t<E>>>)
  constexpr @\exposidnc{error-value-type}@ error_or(G&& def) &&;

  // Deleted: value_or is not available for void expected. Gated to reference E only so
  // that, for value E, no value_or overload is declared at all (there is nothing to
  // delete against).
  template<class U>
    requires is_reference_v<E>
  constexpr void value_or(U&&) const = BEMAN_EXPECTED_DELETE_MSG(
      "expected<void,E>: value_or is not defined for void value_type; there is no "
      "value to fall back from — use has_value()/error()");

  // -------------------------------------------------------------------------
  // [expected.void.monadic] Monadic operations
  // -------------------------------------------------------------------------

  template<class F>
    requires is_constructible_v<E, E&>
  constexpr auto and_then(F&& f) &;
  template<class F>
    requires is_constructible_v<E, E&&>
  constexpr auto and_then(F&& f) &&;
  template<class F>
    requires is_constructible_v<E, const E&>
  constexpr auto and_then(F&& f) const&;
  template<class F>
    requires is_constructible_v<E, const E&&>
  constexpr auto and_then(F&& f) const&&;

  template<class F> constexpr auto or_else(F&& f) &;
  template<class F> constexpr auto or_else(F&& f) &&;
  template<class F> constexpr auto or_else(F&& f) const&;
  template<class F> constexpr auto or_else(F&& f) const&&;

  template<class F>
    requires is_constructible_v<E, E&>
  constexpr auto transform(F&& f) &;
  template<class F>
    requires is_constructible_v<E, E&&>
  constexpr auto transform(F&& f) &&;
  template<class F>
    requires is_constructible_v<E, const E&>
  constexpr auto transform(F&& f) const&;
  template<class F>
    requires is_constructible_v<E, const E&&>
  constexpr auto transform(F&& f) const&&;

  template<class F> constexpr auto transform_error(F&& f) &;
  template<class F> constexpr auto transform_error(F&& f) &&;
  template<class F> constexpr auto transform_error(F&& f) const&;
  template<class F> constexpr auto transform_error(F&& f) const&&;

  // -------------------------------------------------------------------------
  // [expected.void.eq] Equality operators (hidden friends)
  // -------------------------------------------------------------------------

  template<class T2, class E2>
    requires is_void_v<T2>
  friend constexpr bool operator==(const expected& x, const expected<T2, E2>& y);

  template<class E2>
  friend constexpr bool operator==(const expected& x, const unexpected<E2>& e);

private:
  bool @\exposidnc{has-val}@; // exposition only
  union {
    unexpected<E> @\exposidnc{unex}@; // exposition only
  };
};
\end{codeblock}

\begin{itemdescr}
\pnum
\mandates
A program that instantiates the definition of \tcode{expected<T, E>} with an \tcode{E} that is not a valid template argument for \tcode{unexpected} is ill-formed.

\pnum
\remarks
Any object of type \tcode{expected<T, E>} either represents a value of type \tcode{T}, or contains a value of type \tcode{E} nested within it. Member \tcode{has_val} indicates whether the \tcode{expected<T, E>} object represents a value of type \tcode{T}. When \tcode{has_value()} is \tcode{false}, the error is \tcode{unex.error()}.
\end{itemdescr}

\rSec3[expected.void.cons]{Constructors}

\indexlibraryctor{expected}%
\begin{itemdecl}
constexpr expected() noexcept;
\end{itemdecl}

\begin{itemdescr}
\pnum
\ensures
\tcode{has_value()} is \tcode{true}.
\end{itemdescr}

\indexlibraryctor{expected}%
\begin{itemdecl}
constexpr expected(const expected& rhs) noexcept(is_nothrow_copy_constructible_v<E>);
\end{itemdecl}

\begin{itemdescr}
\pnum
\constraints
\tcode{is_copy_constructible_v<E>} is \tcode{true} and \tcode{is_trivially_copy_constructible_v<E>} is \tcode{false}.

\pnum
\effects
If \tcode{rhs.has_value()} is \tcode{false}, direct-non-list-initializes \tcode{unex} with \tcode{rhs.error()}.

\pnum
\ensures
\tcode{rhs.has_value() == this->has_value()}.

\pnum
\throws
Any exception thrown by the initialization of \tcode{unex}.

\pnum
\remarks
This constructor is defined as deleted unless \tcode{is_copy_constructible_v<E>} is \tcode{true} or \tcode{is_reference_v<E>} is \tcode{true}. This constructor is trivial if \tcode{is_trivially_copy_constructible_v<E>} is \tcode{true}.
\end{itemdescr}

\indexlibraryctor{expected}%
\begin{itemdecl}
constexpr expected(expected&& rhs) noexcept(is_nothrow_move_constructible_v<E>);
\end{itemdecl}

\begin{itemdescr}
\pnum
\constraints
\tcode{is_move_constructible_v<E>} is \tcode{true} and \tcode{is_trivially_move_constructible_v<E>} is \tcode{false}.

\pnum
\effects
If \tcode{rhs.has_value()} is \tcode{false}, direct-non-list-initializes \tcode{unex} with \tcode{std::move(rhs.error())}.

\pnum
\ensures
\tcode{rhs.has_value()} is unchanged; \tcode{rhs.has_value() == this->has_value()} is \tcode{true}.

\pnum
\throws
Any exception thrown by the initialization of \tcode{unex}.

\pnum
\remarks
This constructor is trivial if \tcode{is_trivially_move_constructible_v<E>} is \tcode{true}.
\end{itemdescr}

\indexlibraryctor{expected}%
\begin{itemdecl}
template<class U, class G>
  requires(is_void_v<U> && !is_reference_v<E> && !is_same_v<G, E> &&
           is_constructible_v<E, const G&> &&
           !is_constructible_v<unexpected<E>, expected<U, G>&> &&
           !is_constructible_v<unexpected<E>, expected<U, G>&&> &&
           !is_constructible_v<unexpected<E>, const expected<U, G>&> &&
           !is_constructible_v<unexpected<E>, const expected<U, G>&&>)
constexpr explicit(!is_convertible_v<const G&, E>) expected(const expected<U, G>& rhs);
template<class U, class G>
  requires(is_void_v<U> && !is_reference_v<E> && !is_same_v<G, E> &&
           is_constructible_v<E, G> &&
           !is_constructible_v<unexpected<E>, expected<U, G>&> &&
           !is_constructible_v<unexpected<E>, expected<U, G>&&> &&
           !is_constructible_v<unexpected<E>, const expected<U, G>&> &&
           !is_constructible_v<unexpected<E>, const expected<U, G>&&>)
constexpr explicit(!is_convertible_v<G, E>) expected(expected<U, G>&& rhs);
\end{itemdecl}

\begin{itemdescr}
\pnum
\constraints
\tcode{is_void_v<U>} is \tcode{true}; and \tcode{is_constructible_v<E, const G&>} is \tcode{true}; and \tcode{is_constructible_v<unexpected<E>, expected<U, G>&>} is \tcode{false}; and \tcode{is_constructible_v<unexpected<E>, expected<U, G>>} is \tcode{false}; and \tcode{is_constructible_v<unexpected<E>, const expected<U, G>&>} is \tcode{false}; and \tcode{is_constructible_v<unexpected<E>, const expected<U, G>>} is \tcode{false}.

\pnum
\effects
If \tcode{rhs.has_value()} is \tcode{false}, direct-non-list-initializes \tcode{unex} with \tcode{rhs.error()}.

\pnum
\ensures
\tcode{rhs.has_value()} is unchanged; \tcode{rhs.has_value() == this->has_value()} is \tcode{true}.

\pnum
\throws
Any exception thrown by the initialization of \tcode{unex}.
\end{itemdescr}

\indexlibraryctor{expected}%
\begin{itemdecl}
template<class G>
  requires(!is_reference_v<E> && is_constructible_v<E, const G&>)
constexpr explicit(!is_convertible_v<const G&, E>) expected(const unexpected<G>& e);
template<class G>
  requires(!is_reference_v<E> && is_constructible_v<E, G>)
constexpr explicit(!is_convertible_v<G, E>) expected(unexpected<G>&& e);
\end{itemdecl}

\begin{itemdescr}
\pnum
\constraints
\tcode{is_constructible_v<E, const G&>} is \tcode{true}.

\pnum
\effects
Direct-non-list-initializes \tcode{unex} with \tcode{e.error()}.

\pnum
\ensures
\tcode{has_value()} is \tcode{false}.

\pnum
\throws
Any exception thrown by the initialization of \tcode{unex}.
\end{itemdescr}

\indexlibraryctor{expected}%
\begin{itemdecl}
template<class G>
  requires(is_reference_v<E> && is_reference_v<G> && is_constructible_v<E, G> &&
           !reference_constructs_from_temporary_v<E, G>)
constexpr explicit(!is_convertible_v<G, E>) expected(const unexpected<G>& e) noexcept;
template<class G>
  requires(is_reference_v<E> && is_reference_v<G> && is_constructible_v<E, G> &&
           !reference_constructs_from_temporary_v<E, G>)
constexpr explicit(!is_convertible_v<G, E>) expected(unexpected<G>&& e) noexcept;
\end{itemdecl}

\begin{itemdescr}
\pnum
\constraints
\tcode{is_reference_v<G>} is \tcode{true}; and \tcode{is_constructible_v<E, G>} is \tcode{true}; and \tcode{reference_constructs_from_temporary_v<E, G>} is \tcode{false}.

\pnum
\effects
Initializes \tcode{unex} with \tcode{e.error()}.

\pnum
\ensures
\tcode{has_value()} is \tcode{false}.

\pnum
\remarks
This constructor never throws: the referent is bound, not copied.
\end{itemdescr}

\indexlibraryctor{expected}%
\begin{itemdecl}
template<class G>
  requires(is_reference_v<E> && !is_reference_v<G>)
constexpr expected(const unexpected<G>&);
\end{itemdecl}

\begin{itemdescr}
\pnum
\remarks
When \tcode{E} is a reference type, an overload taking \tcode{unexpected<G>} for a non-reference \tcode{G} is defined as deleted: the referent would live inside the (possibly temporary) source \tcode{unexpected<G>} object, and binding \tcode{E&} to it would dangle. Use an \tcode{unexpected<E&>} holding an external object instead.
\end{itemdescr}

\indexlibraryctor{expected}%
\begin{itemdecl}
constexpr explicit expected(in_place_t) noexcept;
\end{itemdecl}

\begin{itemdescr}
\pnum
\ensures
\tcode{has_value()} is \tcode{true}.
\end{itemdescr}

\indexlibraryctor{expected}%
\begin{itemdecl}
template<class... Args>
  requires(is_constructible_v<E, Args...> && !@\exposidnc{unexpect-dangles-v}@<E, Args...>)
constexpr explicit expected(unexpect_t, Args&&... args);
\end{itemdecl}

\begin{itemdescr}
\pnum
\constraints
\tcode{is_constructible_v<E, Args...>} is \tcode{true}.

\pnum
\effects
Direct-non-list-initializes \tcode{unex} with \tcode{std::forward<Args>(args)...}.

\pnum
\ensures
\tcode{has_value()} is \tcode{false}.

\pnum
\throws
Any exception thrown by the initialization of \tcode{unex}.
\end{itemdescr}

\indexlibraryctor{expected}%
\begin{itemdecl}
template<class... Args>
  requires(@\exposidnc{unexpect-dangles-v}@<E, Args...>)
constexpr expected(unexpect_t, Args&&...);
\end{itemdecl}

\begin{itemdescr}
\pnum
\remarks
When \tcode{E} is a reference type, an overload with the same parameter types is defined as deleted if the single argument would bind \tcode{E&} to a temporary, or if it is otherwise not usable to construct \tcode{E}.
\end{itemdescr}

\indexlibraryctor{expected}%
\begin{itemdecl}
template<class U, class... Args>
  requires(!is_reference_v<E> && is_constructible_v<E, initializer_list<U>&, Args...>)
constexpr explicit expected(unexpect_t, initializer_list<U> il, Args&&... args);
\end{itemdecl}

\begin{itemdescr}
\pnum
\constraints
\tcode{is_reference_v<E>} is \tcode{false}, and \tcode{is_constructible_v<E, initializer_list<U>&, Args...>} is \tcode{true}.

\pnum
\effects
Direct-non-list-initializes \tcode{unex} with \tcode{il, std::forward<Args>(args)...}.

\pnum
\ensures
\tcode{has_value()} is \tcode{false}.

\pnum
\throws
Any exception thrown by the initialization of \tcode{unex}.

\pnum
\remarks
An overload with the same parameter types is defined as deleted when \tcode{E} is an lvalue reference type. An initializer list cannot provide the required long-lived error referent.
\end{itemdescr}

\indexlibraryctor{expected}%
\begin{itemdecl}
template<class U, class... Args>
  requires is_reference_v<E>
constexpr expected(unexpect_t, initializer_list<U>, Args&&...);
\end{itemdecl}

\begin{itemdescr}
\pnum
\remarks
An overload with the same parameter types is defined as deleted when \tcode{E} is an lvalue reference type. An initializer list cannot provide the required long-lived error referent.
\end{itemdescr}

\indexlibraryctor{expected}%
\begin{itemdecl}
template<class G>
  requires(is_reference_v<E> && is_convertible_v<G&, E> &&
           !reference_constructs_from_temporary_v<E, G&>)
constexpr explicit(!is_convertible_v<G&, E>) expected(const expected<void, G&>& rhs);
template<class G>
  requires(is_reference_v<E> && is_convertible_v<G&, E> &&
           !reference_constructs_from_temporary_v<E, G&>)
constexpr explicit(!is_convertible_v<G&, E>) expected(expected<void, G&>&& rhs);
\end{itemdecl}

\begin{itemdescr}
\pnum
\constraints
\tcode{is_convertible_v<G&, E>} is \tcode{true}; and \tcode{reference_constructs_from_temporary_v<E, G&>} is \tcode{false}.

\pnum
\effects
If \tcode{rhs.has_value()} is \tcode{false}, direct-non-list-initializes \tcode{unex} with \tcode{rhs.error()}.

\pnum
\ensures
\tcode{rhs.has_value()} is unchanged; \tcode{rhs.has_value() == this->has_value()} is \tcode{true}.

\pnum
\remarks
This constructor never throws: the referent is bound, not copied. It participates in overload resolution only when \tcode{E} is a reference type, mirroring the \tcode{unexpected<G>} reference-\tcode{E} path above.
\end{itemdescr}

\rSec3[expected.void.dtor]{Destructor}

\indexlibrarydtor{expected}%
\begin{itemdecl}
constexpr ~expected();
\end{itemdecl}

\begin{itemdescr}
\pnum
\constraints
\tcode{is_trivially_destructible_v<E>} is \tcode{false}.

\pnum
\effects
If \tcode{has_value()} is \tcode{false}, destroys \tcode{unex}.

\pnum
\remarks
If \tcode{is_trivially_destructible_v<E>} is \tcode{true}, then this destructor is a trivial destructor.
\end{itemdescr}

\rSec3[expected.void.assign]{Assignment}

\indexlibrarymember{operator=}{expected}%
\begin{itemdecl}
constexpr expected& operator=(const expected& rhs) noexcept(
    is_nothrow_copy_constructible_v<E> && is_nothrow_copy_assignable_v<E>);
\end{itemdecl}

\begin{itemdescr}
\pnum
\constraints
\tcode{is_reference_v<E> || (is_copy_constructible_v<E> && is_copy_assignable_v<E>)} is \tcode{true} and \tcode{(is_trivially_copy_constructible_v<E> && is_trivially_copy_assignable_v<E> &&
                           is_trivially_destructible_v<E>)} is \tcode{false}.

\pnum
\effects
If \tcode{this->has_value() && rhs.has_value()}, no effects. Otherwise, if \tcode{this->has_value()}, equivalent to: \tcode{construct_at(addressof(unex), rhs.unex); has_val = false;} Otherwise, if \tcode{rhs.has_value()}, destroys \tcode{unex} and sets \tcode{has_val} to \tcode{true}. Otherwise, equivalent to \tcode{unex = rhs.unex}.

\pnum
\returns
\tcode{*this}.

\pnum
\remarks
This operator is defined as deleted unless \tcode{is_copy_assignable_v<E>} is \tcode{true} or \tcode{is_reference_v<E>} is \tcode{true} and \tcode{is_copy_constructible_v<E>} is \tcode{true} or \tcode{is_reference_v<E>} is \tcode{true}. This operator is trivial if \tcode{is_trivially_copy_constructible_v<E>}, \tcode{is_trivially_copy_assignable_v<E>}, and \tcode{is_trivially_destructible_v<E>} are all \tcode{true}.
\end{itemdescr}

\indexlibrarymember{operator=}{expected}%
\begin{itemdecl}
constexpr expected& operator=(expected&& rhs) noexcept(
    is_nothrow_move_constructible_v<E> && is_nothrow_move_assignable_v<E>);
\end{itemdecl}

\begin{itemdescr}
\pnum
\constraints
\tcode{is_reference_v<E> || (is_move_constructible_v<E> && is_move_assignable_v<E>)} is \tcode{true} and \tcode{(is_trivially_move_constructible_v<E> && is_trivially_move_assignable_v<E> &&
                           is_trivially_destructible_v<E>)} is \tcode{false}.

\pnum
\effects
If \tcode{this->has_value() && rhs.has_value()}, no effects. Otherwise, if \tcode{this->has_value()}, equivalent to: \tcode{construct_at(addressof(unex), std::move(rhs.unex)); has_val = false;} Otherwise, if \tcode{rhs.has_value()}, destroys \tcode{unex} and sets \tcode{has_val} to \tcode{true}. Otherwise, equivalent to \tcode{unex = std::move(rhs.unex)}.

\pnum
\returns
\tcode{*this}.

\pnum
\remarks
The exception specification is equivalent to \tcode{is_nothrow_move_constructible_v<E> && is_nothrow_move_assignable_v<E>}. This operator is trivial if \tcode{is_trivially_move_constructible_v<E>}, \tcode{is_trivially_move_assignable_v<E>}, and \tcode{is_trivially_destructible_v<E>} are all \tcode{true}.
\end{itemdescr}

\indexlibrarymember{operator=}{expected}%
\begin{itemdecl}
template<class G>
  requires(!is_reference_v<E> && is_constructible_v<E, const G&> &&
           is_assignable_v<E&, const G&>)
constexpr expected& operator=(const unexpected<G>& e);
template<class G>
  requires(!is_reference_v<E> && is_constructible_v<E, G> && is_assignable_v<E&, G>)
constexpr expected& operator=(unexpected<G>&& e);
\end{itemdecl}

\begin{itemdescr}
\pnum
\constraints
\tcode{is_constructible_v<E, const G&>} is \tcode{true} and \tcode{is_assignable_v<E&, const G&>} is \tcode{true}.

\pnum
\effects
If \tcode{has_value()} is \tcode{true}, equivalent to: \tcode{construct_at(addressof(unex), e.error()); has_val = false;} Otherwise, equivalent to: \tcode{unex = unexpected<E>(e.error());}

\pnum
\returns
\tcode{*this}.
\end{itemdescr}

\indexlibrarymember{operator=}{expected}%
\begin{itemdecl}
template<class G>
  requires(is_reference_v<E> && is_reference_v<G> && is_constructible_v<E, G> &&
           !reference_constructs_from_temporary_v<E, G>)
constexpr expected& operator=(const unexpected<G>& e);
template<class G>
  requires(is_reference_v<E> && is_reference_v<G> && is_constructible_v<E, G> &&
           !reference_constructs_from_temporary_v<E, G>)
constexpr expected& operator=(unexpected<G>&& e);
\end{itemdecl}

\begin{itemdescr}
\pnum
\constraints
\tcode{is_reference_v<G>} is \tcode{true}; and \tcode{is_constructible_v<E, G>} is \tcode{true}; and \tcode{reference_constructs_from_temporary_v<E, G>} is \tcode{false}.

\pnum
\effects
Rebinds \tcode{unex} to refer to the same object as \tcode{e.error()}.

\pnum
\ensures
\tcode{has_value()} is \tcode{false}.

\pnum
\returns
\tcode{*this}.

\pnum
\remarks
This operator never throws: the referent is bound, not copied.
\end{itemdescr}

\indexlibrarymember{operator=}{expected}%
\begin{itemdecl}
template<class G>
  requires(is_reference_v<E> && !is_reference_v<G>)
constexpr expected& operator=(const unexpected<G>&);
\end{itemdecl}

\begin{itemdescr}
\pnum
\remarks
When \tcode{E} is a reference type, an overload taking \tcode{unexpected<G>} for a non-reference \tcode{G} is defined as deleted: it would rebind \tcode{unex} to \tcode{unexpected<G>}'s temporary storage.
\end{itemdescr}

\indexlibrarymember{emplace}{expected}%
\begin{itemdecl}
constexpr void emplace() noexcept;
\end{itemdecl}

\begin{itemdescr}
\pnum
\effects
If \tcode{has_value()} is \tcode{false}, destroys \tcode{unex} and sets \tcode{has_val} to \tcode{true}.
\end{itemdescr}

\rSec3[expected.void.swap]{Swap}

\indexlibrarymember{swap}{expected}%
\begin{itemdecl}
constexpr void swap(expected& rhs) noexcept(is_nothrow_move_constructible_v<E> &&
                                            (is_reference_v<E> ||
                                             is_nothrow_swappable_v<E>));
\end{itemdecl}

\begin{itemdescr}
\pnum
\constraints
\tcode{is_reference_v<E> || is_swappable_v<E>} is \tcode{true} and \tcode{is_move_constructible_v<E>} is \tcode{true}.

\pnum
\effects
If \tcode{this->has_value()} and \tcode{rhs.has_value()}, no effects. If neither \tcode{*this} nor \tcode{rhs} contains a value, equivalent to \tcode{using std::swap; swap(unex, rhs.unex);}. If \tcode{rhs.has_value()} is \tcode{false} and \tcode{this->has_value()} is \tcode{true}, initializes \tcode{rhs.unex} from \tcode{std::move(unex)}, destroys \tcode{unex}, and leaves \tcode{has_value()} \tcode{false} and \tcode{rhs.has_value()} \tcode{true}. If \tcode{rhs.has_value()} is \tcode{true} and \tcode{this->has_value()} is \tcode{false}, equivalent to \tcode{rhs.swap(*this)}.

\pnum
\throws
Any exception thrown by the expressions in the Effects.

\pnum
\remarks
The exception specification is equivalent to \tcode{is_nothrow_move_constructible_v<E> && (is_reference_v<E> || is_nothrow_swappable_v<E>)}.
\end{itemdescr}

\indexlibraryglobal{swap}%
\begin{itemdecl}
friend constexpr void swap(expected& x, expected& y) noexcept(noexcept(x.swap(y)));
\end{itemdecl}

\begin{itemdescr}
\pnum
\constraints
\tcode{is_reference_v<E> || is_swappable_v<E>} is \tcode{true} and \tcode{is_move_constructible_v<E>} is \tcode{true}.

\pnum
\effects
Equivalent to \tcode{x.swap(y)}.
\end{itemdescr}

\rSec3[expected.void.obs]{Observers}

\indexlibrarymember{operator bool}{expected}%
\indexlibrarymember{has_value}{expected}%
\begin{itemdecl}
constexpr explicit operator bool() const noexcept;
constexpr bool has_value() const noexcept;
\end{itemdecl}

\begin{itemdescr}
\pnum
\returns
\tcode{has_val}.
\end{itemdescr}

\indexlibrarymember{operator*}{expected}%
\begin{itemdecl}
constexpr void operator*() const noexcept;
\end{itemdecl}

\begin{itemdescr}
\pnum
\hardexpects
\tcode{has_value()} is \tcode{true}.
\end{itemdescr}

\indexlibrarymember{value}{expected}%
\begin{itemdecl}
constexpr void value() const&;
\end{itemdecl}

\begin{itemdescr}
\pnum
\mandates
\tcode{is_copy_constructible_v<E>} is \tcode{true}.

\pnum
\throws
\tcode{bad_expected_access(error())} if \tcode{has_value()} is \tcode{false}.
\end{itemdescr}

\indexlibrarymember{value}{expected}%
\begin{itemdecl}
constexpr void value() &&;
\end{itemdecl}

\begin{itemdescr}
\pnum
\mandates
\tcode{is_copy_constructible_v<E>} is \tcode{true} and \tcode{is_move_constructible_v<E>} is \tcode{true}.

\pnum
\throws
\tcode{bad_expected_access(std::move(error()))} if \tcode{has_value()} is \tcode{false}.
\end{itemdescr}

\indexlibrarymember{error}{expected}%
\begin{itemdecl}
constexpr const E& error() const& noexcept;
constexpr E& error() & noexcept;
\end{itemdecl}

\begin{itemdescr}
\pnum
\hardexpects
\tcode{has_value()} is \tcode{false}.

\pnum
\returns
\tcode{unex.error()}.
\end{itemdescr}

\indexlibrarymember{error}{expected}%
\begin{itemdecl}
constexpr const E&& error() const&& noexcept;
constexpr E&& error() && noexcept;
\end{itemdecl}

\begin{itemdescr}
\pnum
\hardexpects
\tcode{has_value()} is \tcode{false}.

\pnum
\returns
\tcode{std::move(unex).error()}.
\end{itemdescr}

\indexlibrarymember{error_or}{expected}%
\begin{itemdecl}
template<class G = @\exposidnc{error-value-type}@>
  requires(is_copy_constructible_v<remove_cv_t<remove_reference_t<E>>> &&
           is_convertible_v<G, remove_cv_t<remove_reference_t<E>>>)
constexpr @\exposidnc{error-value-type}@ error_or(G&& def) const&;
\end{itemdecl}

\begin{itemdescr}
\pnum
\mandates
\tcode{is_copy_constructible_v<error_value_type>} is \tcode{true} and \tcode{is_convertible_v<G, error_value_type>} is \tcode{true}.

\pnum
\returns
\tcode{std::forward<G>(def)} if \tcode{has_value()} is \tcode{true}, \tcode{error()} otherwise.
\end{itemdescr}

\indexlibrarymember{error_or}{expected}%
\begin{itemdecl}
template<class G = @\exposidnc{error-value-type}@>
  requires(is_move_constructible_v<remove_cv_t<remove_reference_t<E>>> &&
           is_convertible_v<G, remove_cv_t<remove_reference_t<E>>>)
constexpr @\exposidnc{error-value-type}@ error_or(G&& def) &&;
\end{itemdecl}

\begin{itemdescr}
\pnum
\mandates
\tcode{is_move_constructible_v<error_value_type>} is \tcode{true} and \tcode{is_convertible_v<G, error_value_type>} is \tcode{true}.

\pnum
\returns
\tcode{std::forward<G>(def)} if \tcode{has_value()} is \tcode{true}, \tcode{std::move(error())} otherwise.
\end{itemdescr}

\indexlibrarymember{value_or}{expected}%
\begin{itemdecl}
template<class U>
  requires is_reference_v<E>
constexpr void value_or(U&&) const;
\end{itemdecl}

\begin{itemdescr}
\pnum
\remarks
\tcode{expected<void, E>} has no \tcode{value_or} member: there is no value to fall back from. This overload exists only to give a clear diagnostic when \tcode{E} is a reference type, and is defined as deleted.
\end{itemdescr}

\rSec3[expected.void.monadic]{Monadic operations}

\indexlibrarymember{and_then}{expected}%
\begin{itemdecl}
template<class F>
  requires is_constructible_v<E, E&>
constexpr auto and_then(F&& f) &;
template<class F>
  requires is_constructible_v<E, const E&>
constexpr auto and_then(F&& f) const&;
\end{itemdecl}

\begin{itemdescr}
\pnum
\constraints
\tcode{is_constructible_v<E, decltype(error())>} is \tcode{true}.

\pnum
\mandates
\tcode{remove_cvref_t<invoke_result_t<F>>} is a specialization of \tcode{expected} and its \tcode{error_type} is the same type as \tcode{E}.

\pnum
\effects
Equivalent to: \tcode{if (has_value()) return invoke(std::forward<F>(f)); else return U(unexpect, error());} where \tcode{U} is \tcode{remove_cvref_t<invoke_result_t<F>>}.
\end{itemdescr}

\indexlibrarymember{and_then}{expected}%
\begin{itemdecl}
template<class F>
  requires is_constructible_v<E, E&&>
constexpr auto and_then(F&& f) &&;
template<class F>
  requires is_constructible_v<E, const E&&>
constexpr auto and_then(F&& f) const&&;
\end{itemdecl}

\begin{itemdescr}
\pnum
\constraints
\tcode{is_constructible_v<E, decltype(std::move(error()))>} is \tcode{true}.

\pnum
\mandates
\tcode{remove_cvref_t<invoke_result_t<F>>} is a specialization of \tcode{expected} and its \tcode{error_type} is the same type as \tcode{E}.

\pnum
\effects
Equivalent to: \tcode{if (has_value()) return invoke(std::forward<F>(f)); else return U(unexpect, std::move(error()));} where \tcode{U} is \tcode{remove_cvref_t<invoke_result_t<F>>}.
\end{itemdescr}

\indexlibrarymember{or_else}{expected}%
\begin{itemdecl}
template<class F> constexpr auto or_else(F&& f) &;
template<class F> constexpr auto or_else(F&& f) const&;
\end{itemdecl}

\begin{itemdescr}
\pnum
\mandates
\tcode{remove_cvref_t<invoke_result_t<F, decltype(error())>>} is a specialization of \tcode{expected} and its \tcode{value_type} is the same type as \tcode{T}.

\pnum
\effects
Equivalent to: \tcode{if (has_value()) return G(); else return invoke(std::forward<F>(f), error());} where \tcode{G} is \tcode{remove_cvref_t<invoke_result_t<F, decltype(error())>>}.
\end{itemdescr}

\indexlibrarymember{or_else}{expected}%
\begin{itemdecl}
template<class F> constexpr auto or_else(F&& f) &&;
template<class F> constexpr auto or_else(F&& f) const&&;
\end{itemdecl}

\begin{itemdescr}
\pnum
\mandates
\tcode{remove_cvref_t<invoke_result_t<F, decltype(std::move(error()))>>} is a specialization of \tcode{expected} and its \tcode{value_type} is the same type as \tcode{T}.

\pnum
\effects
Equivalent to: \tcode{if (has_value()) return G(); else return invoke(std::forward<F>(f), std::move(error()));} where \tcode{G} is \tcode{remove_cvref_t<invoke_result_t<F, decltype(std::move(error()))>>}.
\end{itemdescr}

\indexlibrarymember{transform}{expected}%
\begin{itemdecl}
template<class F>
  requires is_constructible_v<E, E&>
constexpr auto transform(F&& f) &;
template<class F>
  requires is_constructible_v<E, const E&>
constexpr auto transform(F&& f) const&;
\end{itemdecl}

\begin{itemdescr}
\pnum
\constraints
\tcode{is_constructible_v<E, decltype(error())>} is \tcode{true}.

\pnum
\mandates
\tcode{U} is a valid value type for \tcode{expected}, where \tcode{U} is \tcode{remove_cv_t<invoke_result_t<F>>}.

\pnum
\effects
If \tcode{has_value()} is \tcode{false}, returns \tcode{expected<U, E>(unexpect, error())}. Otherwise, if \tcode{is_void_v<U>} is \tcode{false}, returns an \tcode{expected<U, E>} object whose \tcode{has_val} member is \tcode{true} and \tcode{val} member is direct-non-list-initialized with \tcode{invoke(std::forward<F>(f))}. Otherwise, evaluates \tcode{invoke(std::forward<F>(f))} and then returns \tcode{expected<U, E>()}.
\end{itemdescr}

\indexlibrarymember{transform}{expected}%
\begin{itemdecl}
template<class F>
  requires is_constructible_v<E, E&&>
constexpr auto transform(F&& f) &&;
template<class F>
  requires is_constructible_v<E, const E&&>
constexpr auto transform(F&& f) const&&;
\end{itemdecl}

\begin{itemdescr}
\pnum
\constraints
\tcode{is_constructible_v<E, decltype(std::move(error()))>} is \tcode{true}.

\pnum
\mandates
\tcode{U} is a valid value type for \tcode{expected}, where \tcode{U} is \tcode{remove_cv_t<invoke_result_t<F>>}.

\pnum
\effects
If \tcode{has_value()} is \tcode{false}, returns \tcode{expected<U, E>(unexpect, std::move(error()))}. Otherwise, if \tcode{is_void_v<U>} is \tcode{false}, returns an \tcode{expected<U, E>} object whose \tcode{has_val} member is \tcode{true} and \tcode{val} member is direct-non-list-initialized with \tcode{invoke(std::forward<F>(f))}. Otherwise, evaluates \tcode{invoke(std::forward<F>(f))} and then returns \tcode{expected<U, E>()}.
\end{itemdescr}

\indexlibrarymember{transform_error}{expected}%
\begin{itemdecl}
template<class F> constexpr auto transform_error(F&& f) &;
template<class F> constexpr auto transform_error(F&& f) const&;
\end{itemdecl}

\begin{itemdescr}
\pnum
\mandates
\tcode{G} is a valid template argument for \tcode{unexpected} and the declaration \tcode{G g(invoke(std::forward<F>(f), error()));} is well-formed, where \tcode{G} is \tcode{remove_cv_t<invoke_result_t<F, decltype(error())>>}.

\pnum
\returns
If \tcode{has_value()} is \tcode{true}, \tcode{expected<T, G>()}; otherwise, an \tcode{expected<T, G>} object whose \tcode{has_val} member is \tcode{false} and \tcode{unex} member is direct-non-list-initialized with \tcode{invoke(std::forward<F>(f), error())}.
\end{itemdescr}

\indexlibrarymember{transform_error}{expected}%
\begin{itemdecl}
template<class F> constexpr auto transform_error(F&& f) &&;
template<class F> constexpr auto transform_error(F&& f) const&&;
\end{itemdecl}

\begin{itemdescr}
\pnum
\mandates
\tcode{G} is a valid template argument for \tcode{unexpected} and the declaration \tcode{G g(invoke(std::forward<F>(f), std::move(error())));} is well-formed, where \tcode{G} is \tcode{remove_cv_t<invoke_result_t<F, decltype(std::move(error()))>>}.

\pnum
\returns
If \tcode{has_value()} is \tcode{true}, \tcode{expected<T, G>()}; otherwise, an \tcode{expected<T, G>} object whose \tcode{has_val} member is \tcode{false} and \tcode{unex} member is direct-non-list-initialized with \tcode{invoke(std::forward<F>(f), std::move(error()))}.
\end{itemdescr}

\rSec3[expected.void.eq]{Equality operators}

\indexlibraryglobal{operator==}%
\begin{itemdecl}
template<class T2, class E2>
  requires is_void_v<T2>
friend constexpr bool operator==(const expected& x, const expected<T2, E2>& y);
\end{itemdecl}

\begin{itemdescr}
\pnum
\mandates
\tcode{is_void_v<T2>} is \tcode{true}. The expression \tcode{x.error() == y.error()} is well-formed and its result is convertible to \tcode{bool}.

\pnum
\returns
If \tcode{x.has_value() != y.has_value()}, \tcode{false}; otherwise, if \tcode{x.has_value()} is \tcode{true}, \tcode{true}; otherwise \tcode{x.error() == y.error()}.
\end{itemdescr}

\indexlibraryglobal{operator==}%
\begin{itemdecl}
template<class E2>
friend constexpr bool operator==(const expected& x, const unexpected<E2>& e);
\end{itemdecl}

\begin{itemdescr}
\pnum
\mandates
The expression \tcode{x.error() == e.error()} is well-formed and its result is convertible to \tcode{bool}.

\pnum
\returns
\tcode{!x.has_value() && static_cast<bool>(x.error() == e.error())}.
\end{itemdescr}

\rSec2[expected.ref]{Partial specialization of expected for reference types}

\rSec3[expected.ref.general]{General}

\begin{codeblock}
template<class T, class E>
class @\libglobal{expected}@<T&, E> {
private:
  using @\exposidnc{error-value-type}@ = remove_cv_t<remove_reference_t<E>>; // exposition only

public:
  using @\libmember{value_type}{expected}@ = T&;
  using @\libmember{error_type}{expected}@ = E;
  using @\libmember{unexpected_type}{expected}@ = unexpected<E>;

  template<class U> using rebind = expected<U, error_type>;

  // -------------------------------------------------------------------------
  // Constructors
  // -------------------------------------------------------------------------

  expected() = BEMAN_EXPECTED_DELETE_MSG(
      "expected<T&,E>: no default constructor; T& cannot be null");

  // Copy constructor (trivial path). Unconstrained; see the primary
  // template's copy constructor for why.
  constexpr expected(const expected&) = default;

  // Copy constructor (non-trivial path)
  constexpr expected(const expected& rhs) noexcept(is_nothrow_copy_constructible_v<E>)
    requires(is_copy_constructible_v<E> && !is_trivially_copy_constructible_v<E>);

  // Move constructor (trivial path). Unconstrained; no explicit noexcept.
  constexpr expected(expected&&) = default;

  // Move constructor (non-trivial path)
  constexpr expected(expected&& rhs) noexcept(is_nothrow_move_constructible_v<E>)
    requires(is_move_constructible_v<E> && !is_trivially_move_constructible_v<E>);

  // Deleted: no in-place value constructor — T& cannot be constructed in-place
  template<class... Args>
  constexpr expected(in_place_t, Args&&...) = BEMAN_EXPECTED_DELETE_MSG(
      "expected<T&,E>: no in-place value constructor; T& cannot be constructed "
      "in-place — pass a U convertible to T&");

  // Value constructor — takes U that can bind to T&
  template<class U = T>
    requires(!is_same_v<remove_cvref_t<U>, in_place_t> &&
             !is_same_v<remove_cvref_t<U>, expected<T&, E>> &&
             !@\exposidnc{is-unexpected-specialization}@<remove_cvref_t<U>>::value &&
             is_constructible_v<T&, U> && !reference_constructs_from_temporary_v<T&, U>)
  constexpr explicit(!is_convertible_v<U, T&>)
      expected(U&& u) noexcept(is_nothrow_constructible_v<T&, U>);

  // Deleted: binding a temporary to T& creates a dangling reference
  template<class U>
    requires(reference_constructs_from_temporary_v<T&, U>)
  constexpr expected(U&&) =
      BEMAN_EXPECTED_DELETE_MSG("expected<T&,E>: argument would bind a temporary that "
                                "dangles; pass an lvalue reference");

  // Converting constructor from expected<U&, G> (copy) — value-E path
  template<class U, class G>
    requires(!is_reference_v<E> && is_constructible_v<T&, U&> &&
             is_constructible_v<E, const G&> &&
             !reference_constructs_from_temporary_v<T&, U&>)
  constexpr explicit(!is_convertible_v<U&, T&> || !is_convertible_v<const G&, E>)
      expected(const expected<U&, G>& rhs);

  // Converting constructor from expected<U&, G> (move) — value-E path
  template<class U, class G>
    requires(!is_reference_v<E> && is_constructible_v<T&, U&> &&
             is_constructible_v<E, G> && !reference_constructs_from_temporary_v<T&, U&>)
  constexpr explicit(!is_convertible_v<U&, T&> || !is_convertible_v<G, E>)
      expected(expected<U&, G>&& rhs);

  // Converting constructor from expected<U&, G&> (copy/move) — reference-E path: only
  // accepts sources whose error type G is itself a reference convertible to E.
  template<class U, class G>
    requires(is_reference_v<E> && is_reference_v<G> && is_constructible_v<T&, U&> &&
             is_convertible_v<G, E> && !reference_constructs_from_temporary_v<T&, U&>)
  constexpr explicit(!is_convertible_v<U&, T&> || !is_convertible_v<G, E>)
      expected(const expected<U&, G>& rhs);

  template<class U, class G>
    requires(is_reference_v<E> && is_reference_v<G> && is_constructible_v<T&, U&> &&
             is_convertible_v<G, E> && !reference_constructs_from_temporary_v<T&, U&>)
  constexpr explicit(!is_convertible_v<U&, T&> || !is_convertible_v<G, E>)
      expected(expected<U&, G>&& rhs);

  // Constructor from unexpected<G> const& / && — value-E path
  template<class G>
    requires(!is_reference_v<E> && is_constructible_v<E, const G&>)
  constexpr explicit(!is_convertible_v<const G&, E>) expected(const unexpected<G>& e);

  template<class G>
    requires(!is_reference_v<E> && is_constructible_v<E, G>)
  constexpr explicit(!is_convertible_v<G, E>) expected(unexpected<G>&& e);

  // Constructor from unexpected<G> — reference-E path. Allowed only when G is itself a
  // reference, i.e. the source unexpected holds a reference to an external object, so
  // binding E& to e.error() cannot dangle regardless of the source's value category. No
  // const_cast is needed.
  template<class G>
    requires(is_reference_v<E> && is_reference_v<G> && is_constructible_v<E, G> &&
             !reference_constructs_from_temporary_v<E, G>)
  constexpr explicit(!is_convertible_v<G, E>) expected(const unexpected<G>& e) noexcept;

  template<class G>
    requires(is_reference_v<E> && is_reference_v<G> && is_constructible_v<E, G> &&
             !reference_constructs_from_temporary_v<E, G>)
  constexpr explicit(!is_convertible_v<G, E>) expected(unexpected<G>&& e) noexcept;

  // Deleted for reference E with value G: the referent lives inside the unexpected<G>
  // object, so binding E& to it would dangle once a temporary source is destroyed. Use
  // (unexpect, lvalue), or an unexpected<E&> holding an external object, instead.
  template<class G>
    requires(is_reference_v<E> && !is_reference_v<G>)
  constexpr expected(const unexpected<G>&) = BEMAN_EXPECTED_DELETE_MSG(
      "expected<T&,E&>: cannot construct from unexpected<value>; the value would "
      "dangle — use unexpected<E&>");

  template<class G>
    requires(is_reference_v<E> && !is_reference_v<G>)
  constexpr expected(unexpected<G>&&) = BEMAN_EXPECTED_DELETE_MSG(
      "expected<T&,E&>: cannot construct from unexpected<value>; the value would "
      "dangle — use unexpected<E&>");

  // In-place constructor for error
  template<class... Args>
    requires(is_constructible_v<E, Args...> && !@\exposidnc{unexpect-dangles-v}@<E, Args...>)
  constexpr explicit expected(unexpect_t, Args&&... args);

  // Deleted: single argument would bind E& to a temporary — dangling prevention
  template<class... Args>
    requires(@\exposidnc{unexpect-dangles-v}@<E, Args...>)
  constexpr expected(unexpect_t, Args&&...) =
      BEMAN_EXPECTED_DELETE_MSG("expected<T&,E&>: unexpect argument would bind a "
                                "temporary that dangles; pass an lvalue reference");

  // Deleted catch-all: reference E, argument neither constructible nor a dangling case
  template<class... Args>
    requires(is_reference_v<E> && !is_constructible_v<E, Args...> &&
             !@\exposidnc{unexpect-dangles-v}@<E, Args...>)
  constexpr expected(unexpect_t, Args&&...) = BEMAN_EXPECTED_DELETE_MSG(
      "expected<T&,E&>: no viable conversion from the given argument(s) to E&");

  // In-place constructor for error with initializer_list
  template<class U, class... Args>
    requires(!is_reference_v<E> && is_constructible_v<E, initializer_list<U>&, Args...>)
  constexpr explicit expected(unexpect_t, initializer_list<U> il, Args&&... args);

  template<class U, class... Args>
    requires is_reference_v<E>
  constexpr expected(unexpect_t, initializer_list<U>, Args&&...) =
      BEMAN_EXPECTED_DELETE_MSG("expected<T&,E&>: initializer-list error construction "
                                "cannot bind a reference; pass an lvalue reference");

  // -------------------------------------------------------------------------
  // Destructor
  // -------------------------------------------------------------------------

  constexpr ~expected()
    requires is_trivially_destructible_v<E>
  = default;

  constexpr ~expected()
    requires(!is_trivially_destructible_v<E>);

  // -------------------------------------------------------------------------
  // Assignment (rebind semantics)
  // -------------------------------------------------------------------------

  // Copy assignment (trivial path)
  constexpr expected& operator=(const expected&)
    requires(is_trivially_copy_constructible_v<E> &&
             is_trivially_copy_assignable_v<E> && is_trivially_destructible_v<E>)
  = default;

  // Copy assignment (non-trivial path)
  constexpr expected& operator=(const expected& rhs) noexcept(
      is_nothrow_copy_constructible_v<E> && is_nothrow_copy_assignable_v<E>)
    requires((is_reference_v<E> ||
              (is_copy_constructible_v<E> && is_copy_assignable_v<E>)) &&
             !(is_trivially_copy_constructible_v<E> &&
               is_trivially_copy_assignable_v<E> && is_trivially_destructible_v<E>));

  // Move assignment (trivial path)
  constexpr expected& operator=(expected&&) noexcept
    requires(is_trivially_move_constructible_v<E> &&
             is_trivially_move_assignable_v<E> && is_trivially_destructible_v<E>)
  = default;

  // Move assignment (non-trivial path)
  constexpr expected& operator=(expected&& rhs) noexcept(
      is_nothrow_move_constructible_v<E> && is_nothrow_move_assignable_v<E>)
    requires((is_reference_v<E> ||
              (is_move_constructible_v<E> && is_move_assignable_v<E>)) &&
             !(is_trivially_move_constructible_v<E> &&
               is_trivially_move_assignable_v<E> && is_trivially_destructible_v<E>));

  // Rebind reference from lvalue
  template<class U = T>
    requires(!is_same_v<remove_cvref_t<U>, expected<T&, E>> &&
             !@\exposidnc{is-unexpected-specialization}@<remove_cvref_t<U>>::value &&
             is_constructible_v<T&, U> && !reference_constructs_from_temporary_v<T&, U>)
  constexpr expected& operator=(U&& u);

  // Assignment from unexpected<G> — value-E path
  template<class G>
    requires(!is_reference_v<E> && is_constructible_v<E, const G&> &&
             is_assignable_v<E&, const G&>)
  constexpr expected& operator=(const unexpected<G>& e);

  template<class G>
    requires(!is_reference_v<E> && is_constructible_v<E, G> && is_assignable_v<E&, G>)
  constexpr expected& operator=(unexpected<G>&& e);

  // Rebinding assignment for reference E from reference G — binds unex_ to the external
  // referent.
  template<class G>
    requires(is_reference_v<E> && is_reference_v<G> && is_constructible_v<E, G> &&
             !reference_constructs_from_temporary_v<E, G>)
  constexpr expected& operator=(const unexpected<G>& e);

  template<class G>
    requires(is_reference_v<E> && is_reference_v<G> && is_constructible_v<E, G> &&
             !reference_constructs_from_temporary_v<E, G>)
  constexpr expected& operator=(unexpected<G>&& e);

  // Deleted for reference E with value G: would rebind E& to unexpected<G>'s temporary
  // storage.
  template<class G>
    requires(is_reference_v<E> && !is_reference_v<G>)
  constexpr expected& operator=(const unexpected<G>&) = BEMAN_EXPECTED_DELETE_MSG(
      "expected<T&,E&>: cannot assign from unexpected<value>; the value would dangle — "
      "use unexpected<E&>");

  template<class G>
    requires(is_reference_v<E> && !is_reference_v<G>)
  constexpr expected& operator=(unexpected<G>&&) = BEMAN_EXPECTED_DELETE_MSG(
      "expected<T&,E&>: cannot assign from unexpected<value>; the value would dangle — "
      "use unexpected<E&>");

  // emplace — rebind the reference
  template<class U = T>
    requires(is_constructible_v<T&, U> && !reference_constructs_from_temporary_v<T&, U>)
  constexpr T& emplace(U&& u) noexcept(is_nothrow_constructible_v<T&, U>);

  // -------------------------------------------------------------------------
  // Swap
  // -------------------------------------------------------------------------

  constexpr void swap(expected& rhs) noexcept(is_nothrow_move_constructible_v<E> &&
                                              (is_reference_v<E> ||
                                               is_nothrow_swappable_v<E>))
    requires((is_reference_v<E> || is_swappable_v<E>) && is_move_constructible_v<E>);

  friend constexpr void swap(expected& x, expected& y) noexcept(noexcept(x.swap(y)))
    requires((is_reference_v<E> || is_swappable_v<E>) && is_move_constructible_v<E>);

  // -------------------------------------------------------------------------
  // Observers
  // -------------------------------------------------------------------------

  constexpr T* operator->() const noexcept;
  constexpr T& operator*() const noexcept;

  constexpr explicit operator bool() const noexcept;
  constexpr bool has_value() const noexcept;

  constexpr T& value() const&;
  constexpr T& value() &&;

  // error() — shallow const: always returns E& regardless of const on expected
  constexpr const E& error() const& noexcept;
  constexpr E& error() & noexcept;
  constexpr const E&& error() const&& noexcept;
  constexpr E&& error() && noexcept;

  template<class U = remove_cv_t<T>>
    requires(is_object_v<T> && !is_array_v<T>)
  constexpr remove_cv_t<T> value_or(U&& def) const;

  // Constraints spell error_value_type as its underlying trait expression rather than
  // the member typedef: clang (through 22) fails to match an out-of-line constrained
  // member of a partial specialization when the requires-clause names a member typedef
  // of the class.
  template<class G = @\exposidnc{error-value-type}@>
    requires(is_copy_constructible_v<remove_cv_t<remove_reference_t<E>>> &&
             is_convertible_v<G, remove_cv_t<remove_reference_t<E>>>)
  constexpr @\exposidnc{error-value-type}@ error_or(G&& def) const&;

  template<class G = @\exposidnc{error-value-type}@>
    requires(is_move_constructible_v<remove_cv_t<remove_reference_t<E>>> &&
             is_convertible_v<G, remove_cv_t<remove_reference_t<E>>>)
  constexpr @\exposidnc{error-value-type}@ error_or(G&& def) &&;

  // -------------------------------------------------------------------------
  // Monadic operations
  // -------------------------------------------------------------------------

  template<class F>
    requires is_constructible_v<E, E&>
  constexpr auto and_then(F&& f) &;
  template<class F>
    requires is_constructible_v<E, E&&>
  constexpr auto and_then(F&& f) &&;
  template<class F>
    requires is_constructible_v<E, const E&>
  constexpr auto and_then(F&& f) const&;
  template<class F>
    requires is_constructible_v<E, const E&&>
  constexpr auto and_then(F&& f) const&&;

  template<class F> constexpr auto or_else(F&& f) &;
  template<class F> constexpr auto or_else(F&& f) &&;
  template<class F> constexpr auto or_else(F&& f) const&;
  template<class F> constexpr auto or_else(F&& f) const&&;

  // transform: f receives T& (value); error propagates as E; result is expected<U, E>
  template<class F>
    requires is_constructible_v<E, E&>
  constexpr auto transform(F&& f) &;
  template<class F>
    requires is_constructible_v<E, E&&>
  constexpr auto transform(F&& f) &&;
  template<class F>
    requires is_constructible_v<E, const E&>
  constexpr auto transform(F&& f) const&;
  template<class F>
    requires is_constructible_v<E, const E&&>
  constexpr auto transform(F&& f) const&&;

  // transform_error: f receives E; value propagates as T&; result is expected<T&, G>
  template<class F> constexpr auto transform_error(F&& f) &;
  template<class F> constexpr auto transform_error(F&& f) &&;
  template<class F> constexpr auto transform_error(F&& f) const&;
  template<class F> constexpr auto transform_error(F&& f) const&&;

  // -------------------------------------------------------------------------
  // Equality operators (hidden friends)
  // -------------------------------------------------------------------------

  template<class T2, class E2>
    requires(!is_void_v<T2>)
  friend constexpr bool operator==(const expected& x, const expected<T2, E2>& y);

  template<class T2>
    requires(!@\exposidnc{is-expected-specialization}@<T2>::value)
  friend constexpr bool operator==(const expected& x, const T2& val);

  template<class E2>
  friend constexpr bool operator==(const expected& x, const unexpected<E2>& e);

private:
  bool @\exposidnc{has-val}@; // exposition only
  union {
    T* @\exposidnc{val}@;             // exposition only
    unexpected<E> @\exposidnc{unex}@; // exposition only
  };
};
\end{codeblock}

\begin{itemdescr}
\pnum
\mandates
A program that instantiates the definition of \tcode{expected<T&, E>} with an \tcode{E} that is not a valid template argument for \tcode{unexpected} is ill-formed. \tcode{T} shall be an object type that is not an array type.

\pnum
\remarks
An object of type \tcode{expected<T&, E>} either represents a reference to an object of type \tcode{T}, or holds an error. Member \tcode{has_val} indicates whether the object represents a reference. When it represents a reference, member \tcode{val} points to the referenced object, which is not owned by the \tcode{expected} object. Otherwise, the error is \tcode{unex.error()}.
\end{itemdescr}

\rSec3[expected.ref.cons]{Constructors}

\indexlibraryctor{expected}%
\begin{itemdecl}
;
\end{itemdecl}

\begin{itemdescr}
\pnum
\remarks
\tcode{expected<T&, E>} has no default constructor: a reference cannot be null, so there is no empty state to default-construct into.
\end{itemdescr}

\indexlibraryctor{expected}%
\begin{itemdecl}
constexpr expected(const expected&) = default;
\end{itemdecl}

\begin{itemdescr}
\pnum
\effects
If \tcode{rhs.has_value()} is \tcode{true}, initializes \tcode{val} with \tcode{rhs.val}, so that \tcode{*this} and \tcode{rhs} refer to the same object; otherwise, initializes \tcode{unex} with \tcode{rhs.unex}.

\pnum
\ensures
\tcode{rhs.has_value() == this->has_value()}.

\pnum
\remarks
This constructor is trivial.
\end{itemdescr}

\indexlibraryctor{expected}%
\begin{itemdecl}
constexpr expected(const expected& rhs) noexcept(is_nothrow_copy_constructible_v<E>);
constexpr expected(expected&& rhs) noexcept(is_nothrow_move_constructible_v<E>);
\end{itemdecl}

\begin{itemdescr}
\pnum
\constraints
\tcode{is_copy_constructible_v<E>} is \tcode{true} and \tcode{is_trivially_copy_constructible_v<E>} is \tcode{false}.

\pnum
\effects
If \tcode{rhs.has_value()} is \tcode{true}, initializes \tcode{val} with \tcode{rhs.val}, so that \tcode{*this} and \tcode{rhs} refer to the same object; otherwise, initializes \tcode{unex} with \tcode{rhs.unex}.

\pnum
\ensures
\tcode{rhs.has_value() == this->has_value()}.

\pnum
\remarks
This constructor is trivial if the corresponding constructor of \tcode{E} is trivial, and is defined as deleted unless \tcode{is_copy_constructible_v<E>} is \tcode{true}.
\end{itemdescr}

\indexlibraryctor{expected}%
\begin{itemdecl}
constexpr expected(expected&&) = default;
\end{itemdecl}

\begin{itemdescr}
\pnum
\effects
If \tcode{rhs.has_value()} is \tcode{true}, initializes \tcode{val} with \tcode{rhs.val}, so that \tcode{*this} and \tcode{rhs} refer to the same object; otherwise, initializes \tcode{unex} with \tcode{std::move(rhs.unex)}.

\pnum
\ensures
\tcode{rhs.has_value() == this->has_value()}.

\pnum
\remarks
This constructor is trivial.
\end{itemdescr}

\indexlibraryctor{expected}%
\begin{itemdecl}
template<class... Args> constexpr expected(in_place_t, Args&&...);
\end{itemdecl}

\begin{itemdescr}
\pnum
\remarks
\tcode{expected<T&, E>} has no in-place value constructor: \tcode{T&} cannot be constructed in-place. Pass a \tcode{U} convertible to \tcode{T&} instead.
\end{itemdescr}

\indexlibraryctor{expected}%
\begin{itemdecl}
template<class U = T>
  requires(!is_same_v<remove_cvref_t<U>, in_place_t> &&
           !is_same_v<remove_cvref_t<U>, expected<T&, E>> &&
           !@\exposidnc{is-unexpected-specialization}@<remove_cvref_t<U>>::value &&
           is_constructible_v<T&, U> && !reference_constructs_from_temporary_v<T&, U>)
constexpr explicit(!is_convertible_v<U, T&>)
    expected(U&& u) noexcept(is_nothrow_constructible_v<T&, U>);
\end{itemdecl}

\begin{itemdescr}
\pnum
\constraints
\tcode{remove_cvref_t<U>} is not \tcode{in_place_t}, \tcode{expected}, or a specialization of \tcode{unexpected}; \tcode{is_constructible_v<T&, U>} is \tcode{true}; and \tcode{reference_constructs_from_temporary_v<T&, U>} is \tcode{false}.

\pnum
\effects
Let \tcode{r} be the lvalue result of \tcode{T& r = std::forward<U>(u);}. Initializes \tcode{val} with \tcode{addressof(r)}.

\pnum
\ensures
\tcode{has_value()} is \tcode{true}.
\end{itemdescr}

\indexlibraryctor{expected}%
\begin{itemdecl}
template<class U>
  requires(reference_constructs_from_temporary_v<T&, U>)
constexpr expected(U&&);
\end{itemdecl}

\begin{itemdescr}
\pnum
\remarks
A constructor for which \tcode{reference_constructs_from_temporary_v<T&, U>} is \tcode{true} — one that would bind \tcode{T&} to a temporary — is defined as deleted.
\end{itemdescr}

\indexlibraryctor{expected}%
\begin{itemdecl}
template<class U, class G>
  requires(!is_reference_v<E> && is_constructible_v<T&, U&> &&
           is_constructible_v<E, const G&> &&
           !reference_constructs_from_temporary_v<T&, U&>)
constexpr explicit(!is_convertible_v<U&, T&> || !is_convertible_v<const G&, E>)
    expected(const expected<U&, G>& rhs);
template<class U, class G>
  requires(!is_reference_v<E> && is_constructible_v<T&, U&> &&
           is_constructible_v<E, G> && !reference_constructs_from_temporary_v<T&, U&>)
constexpr explicit(!is_convertible_v<U&, T&> || !is_convertible_v<G, E>)
    expected(expected<U&, G>&& rhs);
\end{itemdecl}

\begin{itemdescr}
\pnum
\constraints
\tcode{is_constructible_v<T&, U&>} is \tcode{true}; \tcode{reference_constructs_from_temporary_v<T&, U&>} is \tcode{false}; and \tcode{is_constructible_v<E, const G&>} is \tcode{true}.

\pnum
\effects
If \tcode{rhs.has_value()} is \tcode{true}, initializes \tcode{val} with \tcode{addressof(*rhs)}, so that \tcode{*this} refers to the object referred to by \tcode{rhs}; otherwise, initializes \tcode{unex} with the error of \tcode{rhs}. No object referred to by \tcode{rhs} is moved from.
\end{itemdescr}

\indexlibraryctor{expected}%
\begin{itemdecl}
template<class U, class G>
  requires(is_reference_v<E> && is_reference_v<G> && is_constructible_v<T&, U&> &&
           is_convertible_v<G, E> && !reference_constructs_from_temporary_v<T&, U&>)
constexpr explicit(!is_convertible_v<U&, T&> || !is_convertible_v<G, E>)
    expected(const expected<U&, G>& rhs);
template<class U, class G>
  requires(is_reference_v<E> && is_reference_v<G> && is_constructible_v<T&, U&> &&
           is_convertible_v<G, E> && !reference_constructs_from_temporary_v<T&, U&>)
constexpr explicit(!is_convertible_v<U&, T&> || !is_convertible_v<G, E>)
    expected(expected<U&, G>&& rhs);
\end{itemdecl}

\begin{itemdescr}
\pnum
\constraints
\tcode{is_constructible_v<T&, U&>} is \tcode{true}; \tcode{reference_constructs_from_temporary_v<T&, U&>} is \tcode{false}; \tcode{G} is a reference type; and \tcode{is_convertible_v<G, E>} is \tcode{true}.

\pnum
\effects
If \tcode{rhs.has_value()} is \tcode{true}, initializes \tcode{val} with \tcode{addressof(*rhs)}, so that \tcode{*this} refers to the object referred to by \tcode{rhs}; otherwise, initializes \tcode{unex} with the error of \tcode{rhs}. No object referred to by \tcode{rhs} is moved from.
\end{itemdescr}

\indexlibraryctor{expected}%
\begin{itemdecl}
template<class G>
  requires(!is_reference_v<E> && is_constructible_v<E, const G&>)
constexpr explicit(!is_convertible_v<const G&, E>) expected(const unexpected<G>& e);
template<class G>
  requires(!is_reference_v<E> && is_constructible_v<E, G>)
constexpr explicit(!is_convertible_v<G, E>) expected(unexpected<G>&& e);
\end{itemdecl}

\begin{itemdescr}
\pnum
\constraints
\tcode{is_constructible_v<E, const G&>} is \tcode{true}.

\pnum
\effects
Initializes \tcode{unex} with the error of \tcode{e}.

\pnum
\ensures
\tcode{has_value()} is \tcode{false}.
\end{itemdescr}

\indexlibraryctor{expected}%
\begin{itemdecl}
template<class G>
  requires(is_reference_v<E> && is_reference_v<G> && is_constructible_v<E, G> &&
           !reference_constructs_from_temporary_v<E, G>)
constexpr explicit(!is_convertible_v<G, E>) expected(const unexpected<G>& e) noexcept;
template<class G>
  requires(is_reference_v<E> && is_reference_v<G> && is_constructible_v<E, G> &&
           !reference_constructs_from_temporary_v<E, G>)
constexpr explicit(!is_convertible_v<G, E>) expected(unexpected<G>&& e) noexcept;
\end{itemdecl}

\begin{itemdescr}
\pnum
\constraints
\tcode{is_reference_v<G>} is \tcode{true}; \tcode{is_convertible_v<G, E>} is \tcode{true}; and \tcode{reference_constructs_from_temporary_v<E, G>} is \tcode{false}.

\pnum
\effects
Initializes \tcode{unex} with the error of \tcode{e}.

\pnum
\ensures
\tcode{has_value()} is \tcode{false}.

\pnum
\remarks
This constructor never throws: the referent is bound, not copied.
\end{itemdescr}

\indexlibraryctor{expected}%
\begin{itemdecl}
template<class G>
  requires(is_reference_v<E> && !is_reference_v<G>)
constexpr expected(const unexpected<G>&);
\end{itemdecl}

\begin{itemdescr}
\pnum
\remarks
When \tcode{E} is a reference type, an overload taking \tcode{unexpected<G>} for a non-reference \tcode{G} is defined as deleted: the referent would live inside the (possibly temporary) source \tcode{unexpected<G>} object, and binding \tcode{E&} to it would dangle. Use an \tcode{unexpected<E&>} holding an external object instead.
\end{itemdescr}

\indexlibraryctor{expected}%
\begin{itemdecl}
template<class... Args>
  requires(is_constructible_v<E, Args...> && !@\exposidnc{unexpect-dangles-v}@<E, Args...>)
constexpr explicit expected(unexpect_t, Args&&... args);
\end{itemdecl}

\begin{itemdescr}
\pnum
\constraints
\tcode{is_constructible_v<E, Args...>} is \tcode{true}.

\pnum
\effects
Direct-non-list-initializes \tcode{unex} with \tcode{in_place} and \tcode{std::forward<Args>(args)...}.

\pnum
\ensures
\tcode{has_value()} is \tcode{false}.
\end{itemdescr}

\indexlibraryctor{expected}%
\begin{itemdecl}
template<class... Args>
  requires(@\exposidnc{unexpect-dangles-v}@<E, Args...>)
constexpr expected(unexpect_t, Args&&...);
\end{itemdecl}

\begin{itemdescr}
\pnum
\remarks
When \tcode{E} is a reference type, an overload with the same parameter types is defined as deleted if the single argument would bind \tcode{E&} to a temporary, or if it is otherwise not usable to construct \tcode{E}.
\end{itemdescr}

\indexlibraryctor{expected}%
\begin{itemdecl}
template<class U, class... Args>
  requires(!is_reference_v<E> && is_constructible_v<E, initializer_list<U>&, Args...>)
constexpr explicit expected(unexpect_t, initializer_list<U> il, Args&&... args);
\end{itemdecl}

\begin{itemdescr}
\pnum
\constraints
\tcode{is_reference_v<E>} is \tcode{false}, and \tcode{is_constructible_v<E, initializer_list<U>&, Args...>} is \tcode{true}.

\pnum
\effects
Direct-non-list-initializes \tcode{unex} with \tcode{in_place}, \tcode{il}, and \tcode{std::forward<Args>(args)...}.

\pnum
\ensures
\tcode{has_value()} is \tcode{false}.

\pnum
\remarks
An overload with the same parameter types is defined as deleted when \tcode{E} is an lvalue reference type. An initializer list cannot provide the required long-lived error referent.
\end{itemdescr}

\indexlibraryctor{expected}%
\begin{itemdecl}
template<class U, class... Args>
  requires is_reference_v<E>
constexpr expected(unexpect_t, initializer_list<U>, Args&&...);
\end{itemdecl}

\begin{itemdescr}
\pnum
\remarks
An overload with the same parameter types is defined as deleted when \tcode{E} is an lvalue reference type. An initializer list cannot provide the required long-lived error referent.
\end{itemdescr}

\rSec3[expected.ref.dtor]{Destructor}

\indexlibrarydtor{expected}%
\begin{itemdecl}
constexpr ~expected();
\end{itemdecl}

\begin{itemdescr}
\pnum
\constraints
\tcode{is_trivially_destructible_v<E>} is \tcode{true}.

\pnum
\effects
None: \tcode{*this} never owns the object it refers to; \tcode{T} is never destroyed.

\pnum
\remarks
This destructor is trivial.
\end{itemdescr}

\indexlibrarydtor{expected}%
\begin{itemdecl}
constexpr ~expected();
\end{itemdecl}

\begin{itemdescr}
\pnum
\constraints
\tcode{is_trivially_destructible_v<E>} is \tcode{false}.

\pnum
\effects
If \tcode{has_value()} is \tcode{false}, destroys \tcode{unex}. \tcode{T} is not destroyed; \tcode{*this} never owns the object it refers to.

\pnum
\remarks
This destructor is trivial if \tcode{E} is trivially destructible.
\end{itemdescr}

\rSec3[expected.ref.assign]{Assignment}

\indexlibrarymember{operator=}{expected}%
\begin{itemdecl}
constexpr expected& operator=(const expected&);
\end{itemdecl}

\begin{itemdescr}
\pnum
\constraints
\tcode{is_trivially_copy_constructible_v<E>} is \tcode{true}, \tcode{is_trivially_copy_assignable_v<E>} is \tcode{true}, and \tcode{is_trivially_destructible_v<E>} is \tcode{true}.

\pnum
\effects
If \tcode{rhs.has_value()} is \tcode{true}: if \tcode{has_value()} is \tcode{true}, assigns \tcode{rhs.val} to \tcode{val}; otherwise destroys \tcode{unex} and initializes \tcode{val} with \tcode{rhs.val}. If \tcode{rhs.has_value()} is \tcode{false}, the error of \tcode{rhs} is assigned to or used to initialize \tcode{unex}, as for the primary template. In every case \tcode{*this} comes to refer to the object \tcode{rhs} refers to, or to hold the error of \tcode{rhs}.

\pnum
\returns
\tcode{*this}.

\pnum
\remarks
Assignment rebinds: assigning to an \tcode{expected<T&, E>} that holds a value changes which object it refers to. It never assigns through to the referent. This operator is trivial.
\end{itemdescr}

\indexlibrarymember{operator=}{expected}%
\begin{itemdecl}
constexpr expected& operator=(const expected& rhs) noexcept(
    is_nothrow_copy_constructible_v<E> && is_nothrow_copy_assignable_v<E>);
constexpr expected& operator=(expected&& rhs) noexcept(
    is_nothrow_move_constructible_v<E> && is_nothrow_move_assignable_v<E>);
\end{itemdecl}

\begin{itemdescr}
\pnum
\constraints
\tcode{is_reference_v<E> || (is_copy_constructible_v<E> && is_copy_assignable_v<E>)} is \tcode{true} and \tcode{(is_trivially_copy_constructible_v<E> && is_trivially_copy_assignable_v<E> &&
                           is_trivially_destructible_v<E>)} is \tcode{false}.

\pnum
\effects
If \tcode{rhs.has_value()} is \tcode{true}: if \tcode{has_value()} is \tcode{true}, assigns \tcode{rhs.val} to \tcode{val}; otherwise destroys \tcode{unex} and initializes \tcode{val} with \tcode{rhs.val}. If \tcode{rhs.has_value()} is \tcode{false}, the error of \tcode{rhs} is assigned to or used to initialize \tcode{unex}, as for the primary template. In every case \tcode{*this} comes to refer to the object \tcode{rhs} refers to, or to hold the error of \tcode{rhs}.

\pnum
\returns
\tcode{*this}.

\pnum
\remarks
This operator is defined as deleted unless \tcode{is_copy_assignable_v<E>} is \tcode{true} and \tcode{is_copy_constructible_v<E>} is \tcode{true}.
\end{itemdescr}

\indexlibrarymember{operator=}{expected}%
\begin{itemdecl}
constexpr expected& operator=(expected&&) noexcept;
\end{itemdecl}

\begin{itemdescr}
\pnum
\constraints
\tcode{is_trivially_move_constructible_v<E>} is \tcode{true}, \tcode{is_trivially_move_assignable_v<E>} is \tcode{true}, and \tcode{is_trivially_destructible_v<E>} is \tcode{true}.

\pnum
\effects
If \tcode{rhs.has_value()} is \tcode{true}: if \tcode{has_value()} is \tcode{true}, assigns \tcode{rhs.val} to \tcode{val}; otherwise destroys \tcode{unex} and initializes \tcode{val} with \tcode{rhs.val}. If \tcode{rhs.has_value()} is \tcode{false}, the error of \tcode{rhs} is assigned to or used to initialize \tcode{unex}, as for the primary template.

\pnum
\returns
\tcode{*this}.

\pnum
\remarks
This operator is trivial.
\end{itemdescr}

\indexlibrarymember{operator=}{expected}%
\begin{itemdecl}
template<class U = T>
  requires(!is_same_v<remove_cvref_t<U>, expected<T&, E>> &&
           !@\exposidnc{is-unexpected-specialization}@<remove_cvref_t<U>>::value &&
           is_constructible_v<T&, U> && !reference_constructs_from_temporary_v<T&, U>)
constexpr expected& operator=(U&& u);
\end{itemdecl}

\begin{itemdescr}
\pnum
\constraints
\tcode{remove_cvref_t<U>} is neither \tcode{expected} nor a specialization of \tcode{unexpected}, \tcode{is_constructible_v<T&, U>} is \tcode{true}, and \tcode{reference_constructs_from_temporary_v<T&, U>} is \tcode{false}.

\pnum
\effects
Let \tcode{r} be the lvalue result of \tcode{T& r = std::forward<U>(u);}. If \tcode{has_value()} is \tcode{true}, assigns \tcode{addressof(r)} to \tcode{val}. Otherwise, destroys \tcode{unex}, initializes \tcode{val} with \tcode{addressof(r)}, and sets \tcode{has_val} to \tcode{true}; if binding \tcode{r} throws, \tcode{*this} is left unchanged.

\pnum
\returns
\tcode{*this}.
\end{itemdescr}

\indexlibrarymember{operator=}{expected}%
\begin{itemdecl}
template<class G>
  requires(!is_reference_v<E> && is_constructible_v<E, const G&> &&
           is_assignable_v<E&, const G&>)
constexpr expected& operator=(const unexpected<G>& e);
template<class G>
  requires(!is_reference_v<E> && is_constructible_v<E, G> && is_assignable_v<E&, G>)
constexpr expected& operator=(unexpected<G>&& e);
\end{itemdecl}

\begin{itemdescr}
\pnum
\constraints
\tcode{is_constructible_v<E, const G&>} is \tcode{true} and \tcode{is_assignable_v<E&, const G&>} is \tcode{true}.

\pnum
\effects
Makes \tcode{*this} hold the error of \tcode{e}, reinitializing \tcode{unex} from \tcode{e} rather than assigning through it.

\pnum
\returns
\tcode{*this}.
\end{itemdescr}

\indexlibrarymember{operator=}{expected}%
\begin{itemdecl}
template<class G>
  requires(is_reference_v<E> && is_reference_v<G> && is_constructible_v<E, G> &&
           !reference_constructs_from_temporary_v<E, G>)
constexpr expected& operator=(const unexpected<G>& e);
template<class G>
  requires(is_reference_v<E> && is_reference_v<G> && is_constructible_v<E, G> &&
           !reference_constructs_from_temporary_v<E, G>)
constexpr expected& operator=(unexpected<G>&& e);
\end{itemdecl}

\begin{itemdescr}
\pnum
\constraints
\tcode{is_reference_v<G>} is \tcode{true}; \tcode{is_convertible_v<G, E>} is \tcode{true}; and \tcode{reference_constructs_from_temporary_v<E, G>} is \tcode{false}.

\pnum
\effects
Makes \tcode{*this} hold the error of \tcode{e}, reinitializing \tcode{unex} from \tcode{e} rather than assigning through it. \tcode{unex.error()} thereafter refers to the same object as \tcode{e.error()}; the previously referenced object, if any, is not modified.

\pnum
\returns
\tcode{*this}.

\pnum
\remarks
This operator never throws: the referent is bound, not copied.
\end{itemdescr}

\indexlibrarymember{operator=}{expected}%
\begin{itemdecl}
template<class G>
  requires(is_reference_v<E> && !is_reference_v<G>)
constexpr expected& operator=(const unexpected<G>&);
\end{itemdecl}

\begin{itemdescr}
\pnum
\remarks
When \tcode{E} is a reference type, an overload taking \tcode{unexpected<G>} for a non-reference \tcode{G} is defined as deleted: it would rebind \tcode{unex} to \tcode{unexpected<G>}'s temporary storage.
\end{itemdescr}

\indexlibrarymember{emplace}{expected}%
\begin{itemdecl}
template<class U = T>
  requires(is_constructible_v<T&, U> && !reference_constructs_from_temporary_v<T&, U>)
constexpr T& emplace(U&& u) noexcept(is_nothrow_constructible_v<T&, U>);
\end{itemdecl}

\begin{itemdescr}
\pnum
\constraints
\tcode{is_constructible_v<T&, U>} is \tcode{true} and \tcode{reference_constructs_from_temporary_v<T&, U>} is \tcode{false}.

\pnum
\effects
Rebinds \tcode{*this} to refer to the object bound by \tcode{T& r = std::forward<U>(u);}: if \tcode{has_value()} is \tcode{false}, destroys \tcode{unex} first. Sets \tcode{val} to \tcode{addressof(r)} and \tcode{has_val} to \tcode{true}.

\pnum
\returns
\tcode{*val}.
\end{itemdescr}

\rSec3[expected.ref.swap]{Swap}

\indexlibrarymember{swap}{expected}%
\begin{itemdecl}
constexpr void swap(expected& rhs) noexcept(is_nothrow_move_constructible_v<E> &&
                                            (is_reference_v<E> ||
                                             is_nothrow_swappable_v<E>));
\end{itemdecl}

\begin{itemdescr}
\pnum
\constraints
\tcode{is_reference_v<E> || is_swappable_v<E>} is \tcode{true} and \tcode{is_move_constructible_v<E>} is \tcode{true}.

\pnum
\effects
Exchanges the states of \tcode{*this} and \tcode{rhs}. When both hold values, exchanges \tcode{val} and \tcode{rhs.val} — the referenced objects are not swapped. Otherwise behaves as the primary template's \tcode{swap} does for the error.

\pnum
\remarks
The exception specification is equivalent to \tcode{is_nothrow_move_constructible_v<E> && (is_reference_v<E> || is_nothrow_swappable_v<E>)}.
\end{itemdescr}

\indexlibraryglobal{swap}%
\begin{itemdecl}
friend constexpr void swap(expected& x, expected& y) noexcept(noexcept(x.swap(y)));
\end{itemdecl}

\begin{itemdescr}
\pnum
\constraints
\tcode{is_reference_v<E> || is_swappable_v<E>} is \tcode{true} and \tcode{is_move_constructible_v<E>} is \tcode{true}.

\pnum
\effects
Equivalent to \tcode{x.swap(y)}.
\end{itemdescr}

\rSec3[expected.ref.obs]{Observers}

\indexlibrarymember{operator->}{expected}%
\begin{itemdecl}
constexpr T* operator->() const noexcept;
\end{itemdecl}

\begin{itemdescr}
\pnum
\expects
\tcode{has_value()} is \tcode{true}.

\pnum
\returns
\tcode{val}.

\pnum
\remarks
This is a \tcode{const} member function that returns a non-\tcode{const} \tcode{T*}; the constness of \tcode{*this} does not propagate to the referenced object. For deep \tcode{const}, use \tcode{expected<const T&, E>}.
\end{itemdescr}

\indexlibrarymember{operator*}{expected}%
\begin{itemdecl}
constexpr T& operator*() const noexcept;
\end{itemdecl}

\begin{itemdescr}
\pnum
\expects
\tcode{has_value()} is \tcode{true}.

\pnum
\returns
\tcode{*val}.

\pnum
\remarks
This is a \tcode{const} member function that returns a non-\tcode{const} \tcode{T&}; the constness of \tcode{*this} does not propagate to the referenced object. For deep \tcode{const}, use \tcode{expected<const T&, E>}.
\end{itemdescr}

\indexlibrarymember{operator bool}{expected}%
\indexlibrarymember{has_value}{expected}%
\begin{itemdecl}
constexpr explicit operator bool() const noexcept;
constexpr bool has_value() const noexcept;
\end{itemdecl}

\begin{itemdescr}
\pnum
\returns
\tcode{has_val}.
\end{itemdescr}

\indexlibrarymember{value}{expected}%
\begin{itemdecl}
constexpr T& value() const&;
\end{itemdecl}

\begin{itemdescr}
\pnum
\mandates
\tcode{is_copy_constructible_v<\exposid{error-value-type}>} is \tcode{true}.

\pnum
\returns
\tcode{*val} if \tcode{has_value()} is \tcode{true}.

\pnum
\throws
\tcode{bad_expected_access(as_const(error()))} if \tcode{has_value()} is \tcode{false}.
\end{itemdescr}

\indexlibrarymember{value}{expected}%
\begin{itemdecl}
constexpr T& value() &&;
\end{itemdecl}

\begin{itemdescr}
\pnum
\returns
\tcode{*val} if \tcode{has_value()} is \tcode{true}.

\pnum
\throws
\tcode{bad_expected_access(std::move(error()))} if \tcode{has_value()} is \tcode{false}.
\end{itemdescr}

\indexlibrarymember{error}{expected}%
\begin{itemdecl}
constexpr const E& error() const& noexcept;
constexpr E& error() & noexcept;
\end{itemdecl}

\begin{itemdescr}
\pnum
\expects
\tcode{has_value()} is \tcode{false}.

\pnum
\returns
\tcode{unex.error()}.
\end{itemdescr}

\indexlibrarymember{error}{expected}%
\begin{itemdecl}
constexpr const E&& error() const&& noexcept;
constexpr E&& error() && noexcept;
\end{itemdecl}

\begin{itemdescr}
\pnum
\expects
\tcode{has_value()} is \tcode{false}.

\pnum
\returns
\tcode{std::move(unex).error()}.
\end{itemdescr}

\indexlibrarymember{value_or}{expected}%
\begin{itemdecl}
template<class U = remove_cv_t<T>>
  requires(is_object_v<T> && !is_array_v<T>)
constexpr remove_cv_t<T> value_or(U&& def) const;
\end{itemdecl}

\begin{itemdescr}
\pnum
\mandates
\tcode{is_convertible_v<T&, remove_cv_t<T>>} and \tcode{is_convertible_v<U, remove_cv_t<T>>} are \tcode{true}.

\pnum
\returns
\tcode{has_value() ? static_cast<remove_cv_t<T>>(*val) : static_cast<remove_cv_t<T>>(std::forward<U>(def))}. The result is an object, never a reference.
\end{itemdescr}

\indexlibrarymember{error_or}{expected}%
\begin{itemdecl}
template<class G = @\exposidnc{error-value-type}@>
  requires(is_copy_constructible_v<remove_cv_t<remove_reference_t<E>>> &&
           is_convertible_v<G, remove_cv_t<remove_reference_t<E>>>)
constexpr @\exposidnc{error-value-type}@ error_or(G&& def) const&;
\end{itemdecl}

\begin{itemdescr}
\pnum
\returns
\tcode{std::forward<G>(def)} if \tcode{has_value()} is \tcode{true}, \tcode{error()} otherwise. The result is an object, never a reference.
\end{itemdescr}

\indexlibrarymember{error_or}{expected}%
\begin{itemdecl}
template<class G = @\exposidnc{error-value-type}@>
  requires(is_move_constructible_v<remove_cv_t<remove_reference_t<E>>> &&
           is_convertible_v<G, remove_cv_t<remove_reference_t<E>>>)
constexpr @\exposidnc{error-value-type}@ error_or(G&& def) &&;
\end{itemdecl}

\begin{itemdescr}
\pnum
\returns
\tcode{std::forward<G>(def)} if \tcode{has_value()} is \tcode{true}, \tcode{std::move(error())} otherwise. The result is an object, never a reference.
\end{itemdescr}

\rSec3[expected.ref.monadic]{Monadic operations}

\indexlibrarymember{and_then}{expected}%
\begin{itemdecl}
template<class F>
  requires is_constructible_v<E, E&>
constexpr auto and_then(F&& f) &;
template<class F>
  requires is_constructible_v<E, const E&>
constexpr auto and_then(F&& f) const&;
\end{itemdecl}

\begin{itemdescr}
\pnum
\constraints
\tcode{is_constructible_v<E, decltype(error())>} is \tcode{true}.

\pnum
\mandates
\tcode{remove_cvref_t<invoke_result_t<F, T&>>} is a specialization of \tcode{expected} and its \tcode{error_type} is the same type as \tcode{E}.

\pnum
\effects
Equivalent to: \tcode{if (has_value()) return invoke(std::forward<F>(f), *val); else return U(unexpect, error());} where \tcode{U} is \tcode{remove_cvref_t<invoke_result_t<F, T&>>}.

\pnum
\remarks
The member templates \tcode{and_then}, \tcode{or_else}, \tcode{transform}, and \tcode{transform_error} behave as specified for the primary template, with one difference: the value is passed to the callable as \tcode{T&} for every ref-qualification of \tcode{*this}. An rvalue \tcode{expected<T&, E>} does not pass its referent as an rvalue; the object referred to is never moved from by these operations.
\end{itemdescr}

\indexlibrarymember{and_then}{expected}%
\begin{itemdecl}
template<class F>
  requires is_constructible_v<E, E&&>
constexpr auto and_then(F&& f) &&;
template<class F>
  requires is_constructible_v<E, const E&&>
constexpr auto and_then(F&& f) const&&;
\end{itemdecl}

\begin{itemdescr}
\pnum
\constraints
\tcode{is_constructible_v<E, decltype(std::move(error()))>} is \tcode{true}.

\pnum
\mandates
\tcode{remove_cvref_t<invoke_result_t<F, T&>>} is a specialization of \tcode{expected} and its \tcode{error_type} is the same type as \tcode{E}.

\pnum
\effects
Equivalent to: \tcode{if (has_value()) return invoke(std::forward<F>(f), *val); else return U(unexpect, std::move(error()));} where \tcode{U} is \tcode{remove_cvref_t<invoke_result_t<F, T&>>}.
\end{itemdescr}

\indexlibrarymember{or_else}{expected}%
\begin{itemdecl}
template<class F> constexpr auto or_else(F&& f) &;
template<class F> constexpr auto or_else(F&& f) const&;
\end{itemdecl}

\begin{itemdescr}
\pnum
\mandates
\tcode{remove_cvref_t<invoke_result_t<F, decltype(error())>>} is a specialization of \tcode{expected} and its \tcode{value_type} is the same type as \tcode{T&}.

\pnum
\effects
Equivalent to: \tcode{if (has_value()) return G(*val); else return invoke(std::forward<F>(f), error());} where \tcode{G} is \tcode{remove_cvref_t<invoke_result_t<F, decltype(error())>>}.
\end{itemdescr}

\indexlibrarymember{or_else}{expected}%
\begin{itemdecl}
template<class F> constexpr auto or_else(F&& f) &&;
template<class F> constexpr auto or_else(F&& f) const&&;
\end{itemdecl}

\begin{itemdescr}
\pnum
\mandates
\tcode{remove_cvref_t<invoke_result_t<F, decltype(std::move(error()))>>} is a specialization of \tcode{expected} and its \tcode{value_type} is the same type as \tcode{T&}.

\pnum
\effects
Equivalent to: \tcode{if (has_value()) return G(*val); else return invoke(std::forward<F>(f), std::move(error()));} where \tcode{G} is \tcode{remove_cvref_t<invoke_result_t<F, decltype(std::move(error()))>>}.
\end{itemdescr}

\indexlibrarymember{transform}{expected}%
\begin{itemdecl}
template<class F>
  requires is_constructible_v<E, E&>
constexpr auto transform(F&& f) &;
template<class F>
  requires is_constructible_v<E, const E&>
constexpr auto transform(F&& f) const&;
\end{itemdecl}

\begin{itemdescr}
\pnum
\constraints
\tcode{is_constructible_v<E, decltype(error())>} is \tcode{true}.

\pnum
\effects
Equivalent to: \tcode{if (!has_value()) return U(unexpect, error()); else return expected<U2, E>(in_place, invoke(std::forward<F>(f), *val));} where \tcode{U2} is \tcode{remove_cv_t<invoke_result_t<F, T&>>} and \tcode{U} is \tcode{expected<U2, E>}.
\end{itemdescr}

\indexlibrarymember{transform}{expected}%
\begin{itemdecl}
template<class F>
  requires is_constructible_v<E, E&&>
constexpr auto transform(F&& f) &&;
template<class F>
  requires is_constructible_v<E, const E&&>
constexpr auto transform(F&& f) const&&;
\end{itemdecl}

\begin{itemdescr}
\pnum
\constraints
\tcode{is_constructible_v<E, decltype(std::move(error()))>} is \tcode{true}.

\pnum
\effects
Equivalent to: \tcode{if (!has_value()) return U(unexpect, std::move(error())); else return expected<U2, E>(in_place, invoke(std::forward<F>(f), *val));} where \tcode{U2} is \tcode{remove_cv_t<invoke_result_t<F, T&>>} and \tcode{U} is \tcode{expected<U2, E>}.
\end{itemdescr}

\indexlibrarymember{transform_error}{expected}%
\begin{itemdecl}
template<class F> constexpr auto transform_error(F&& f) &;
template<class F> constexpr auto transform_error(F&& f) const&;
\end{itemdecl}

\begin{itemdescr}
\pnum
\mandates
\begin{itemize}
\item \tcode{is_object_v<G>} is \tcode{true},
\item \tcode{is_array_v<G>} is \tcode{false},
\item \tcode{is_same_v<G, remove_cv_t<G>>} is \tcode{true},
\item \tcode{\exposid{is-unexpected-specialization}<G>::value} is \tcode{false}.
\end{itemize}

\pnum
\effects
Equivalent to: \tcode{if (has_value()) return G(*val); else return expected<T&, G2>(unexpect, invoke(std::forward<F>(f), error()));} where \tcode{G2} is \tcode{remove_cv_t<invoke_result_t<F, decltype(error())>>} and \tcode{G} is \tcode{expected<T&, G2>}.
\end{itemdescr}

\indexlibrarymember{transform_error}{expected}%
\begin{itemdecl}
template<class F> constexpr auto transform_error(F&& f) &&;
template<class F> constexpr auto transform_error(F&& f) const&&;
\end{itemdecl}

\begin{itemdescr}
\pnum
\mandates
\begin{itemize}
\item \tcode{is_object_v<G>} is \tcode{true},
\item \tcode{is_array_v<G>} is \tcode{false},
\item \tcode{is_same_v<G, remove_cv_t<G>>} is \tcode{true},
\item \tcode{\exposid{is-unexpected-specialization}<G>::value} is \tcode{false}.
\end{itemize}

\pnum
\effects
Equivalent to: \tcode{if (has_value()) return G(*val); else return expected<T&, G2>(unexpect, invoke(std::forward<F>(f), std::move(error())));} where \tcode{G2} is \tcode{remove_cv_t<invoke_result_t<F, decltype(std::move(error()))>>} and \tcode{G} is \tcode{expected<T&, G2>}.
\end{itemdescr}

\rSec3[expected.ref.eq]{Equality operators}

\indexlibraryglobal{operator==}%
\begin{itemdecl}
template<class T2, class E2>
  requires(!is_void_v<T2>)
friend constexpr bool operator==(const expected& x, const expected<T2, E2>& y);
\end{itemdecl}

\begin{itemdescr}
\pnum
\mandates
\tcode{!is_void_v<T2>} is \tcode{true}. The expression \tcode{*x == *y} is well-formed and its result is convertible to \tcode{bool}. The expression \tcode{x.error() == y.error()} is well-formed and its result is convertible to \tcode{bool}.

\pnum
\returns
If \tcode{x.has_value() != y.has_value()}, \tcode{false}; otherwise, if \tcode{x.has_value()} is \tcode{true}, \tcode{*x == *y}; otherwise \tcode{x.error() == y.error()}.

\pnum
\remarks
The equality operators behave as specified for the primary template, comparing referents through \tcode{operator*} and errors through \tcode{error()}.
\end{itemdescr}

\indexlibraryglobal{operator==}%
\begin{itemdecl}
template<class T2>
  requires(!@\exposidnc{is-expected-specialization}@<T2>::value)
friend constexpr bool operator==(const expected& x, const T2& val);
\end{itemdecl}

\begin{itemdescr}
\pnum
\mandates
\tcode{T2} is not a specialization of \tcode{expected}. The expression \tcode{*x == val} is well-formed and its result is convertible to \tcode{bool}.

\pnum
\returns
\tcode{x.has_value() && static_cast<bool>(*x == val)}.
\end{itemdescr}

\indexlibraryglobal{operator==}%
\begin{itemdecl}
template<class E2>
friend constexpr bool operator==(const expected& x, const unexpected<E2>& e);
\end{itemdecl}

\begin{itemdescr}
\pnum
\mandates
The expression \tcode{x.error() == e.error()} is well-formed and its result is convertible to \tcode{bool}.

\pnum
\returns
\tcode{!x.has_value() && static_cast<bool>(x.error() == e.error())}.
\end{itemdescr}
